\documentclass[12pt,a4paper]{article}

% --- Metadata (per course) ---
\def\courseshortheader{XÁC SUẤT}
\def\coverbookline{XÁC SUẤT}
\def\covermaintitle{Xác suất}
\def\coversubtitle{Giáo án lý thuyết - Biến cố, Biến ngẫu nhiên, phân phối, CLT}

% Shared preamble for nhom_5 (requires \courseshortheader before \input)
\usepackage[utf8]{inputenc}
\usepackage[vietnamese]{babel}
\usepackage{amsmath,amssymb,amsthm}
\usepackage{geometry}
\usepackage{enumitem}
\usepackage[most]{tcolorbox}
\usepackage{hyperref}
\usepackage{fancyhdr}
\usepackage{graphicx}
\usepackage{booktabs}
\usepackage{tikz}
\usetikzlibrary{calc,arrows.meta,decorations.pathreplacing,positioning}
\usepackage{eso-pic}
\usepackage{xcolor}

\geometry{a4paper, margin=2cm, bottom=2.5cm}
\hypersetup{colorlinks=true, linkcolor=blue!70!black, urlcolor=blue}
\setlist[itemize]{topsep=4pt, itemsep=2pt}
\setlist[enumerate]{topsep=4pt, itemsep=2pt}

\AddToShipoutPictureFG{%
  \AtPageCenter{%
    \begin{tikzpicture}[remember picture, overlay]
      \node[rotate=45, scale=1, opacity=0.06]{%
        \fontsize{60}{72}\selectfont\bfseries\sffamily Đặng Tấn Quí%
      };
    \end{tikzpicture}%
  }%
}

\pagestyle{fancy}
\fancyhf{}
\fancyhead[L]{\small\textit{\courseshortheader}}
\fancyhead[R]{\small\textit{Đặng Tấn Quí}}
\fancyfoot[C]{\thepage}
\fancyfoot[L]{\footnotesize\textcopyright\ Đặng Tấn Quí}
\fancyfoot[R]{\footnotesize SĐT: 0377\,122\,210}
\renewcommand{\headrulewidth}{0.4pt}
\renewcommand{\footrulewidth}{0.4pt}

\fancypagestyle{plain}{%
  \fancyhf{}
  \fancyfoot[C]{\thepage}
  \fancyfoot[L]{\footnotesize\textcopyright\ Đặng Tấn Quí}
  \fancyfoot[R]{\footnotesize SĐT: 0377\,122\,210}
  \renewcommand{\headrulewidth}{0pt}
  \renewcommand{\footrulewidth}{0.4pt}
}

\newtcolorbox{dongco}[1][]{%
  enhanced, breakable,
  colback=teal!4!white, colframe=teal!60!black,
  fonttitle=\bfseries, title={Động cơ học -- Tại sao cần khái niệm này?}, #1%
}
\newtcolorbox{ynghia}[1][]{%
  enhanced, breakable,
  colback=purple!3!white, colframe=purple!50!black,
  fonttitle=\bfseries, title={Ý nghĩa \& Trực giác}, #1%
}
\newtcolorbox{ungdung}[1][]{%
  enhanced, breakable,
  colback=olive!5!white, colframe=olive!60!black,
  fonttitle=\bfseries, title={Ứng dụng thực tế}, #1%
}
\newtcolorbox{dinhnghia}[1][]{%
  enhanced, breakable,
  colback=blue!3!white, colframe=blue!60!black,
  fonttitle=\bfseries, title={Định nghĩa}, #1%
}
\newtcolorbox{dinhly}[1][]{%
  enhanced, breakable,
  colback=green!3!white, colframe=green!50!black,
  fonttitle=\bfseries, title={Định lý}, #1%
}
\newtcolorbox{tinhchat}[1][]{%
  enhanced, breakable,
  colback=yellow!5!white, colframe=orange!50!black,
  fonttitle=\bfseries, title={Tính chất}, #1%
}
\newtcolorbox{phuongphap}[1][]{%
  enhanced, breakable,
  colback=gray!5!white, colframe=gray!50!black,
  fonttitle=\bfseries, title={Phương pháp}, #1%
}
\newtcolorbox{luuy}[1][]{%
  enhanced, breakable,
  colback=red!3!white, colframe=red!50!black,
  fonttitle=\bfseries, title={Lưu ý}, #1%
}
\newtcolorbox{congthuc}[1][]{%
  enhanced, breakable,
  colback=cyan!3!white, colframe=cyan!50!black,
  fonttitle=\bfseries, title={Công thức}, #1%
}
\newtcolorbox{vidu}[1][]{%
  enhanced, breakable,
  colback=violet!3!white, colframe=violet!50!black,
  fonttitle=\bfseries, title={Ví dụ}, #1%
}
\newtcolorbox{phanvidu}[1][]{%
  enhanced, breakable,
  colback=red!2!white, colframe=red!40!black,
  fonttitle=\bfseries, title={Phản ví dụ}, #1%
}
\newtcolorbox{tongket}[1][]{%
  enhanced, breakable,
  colback=blue!4!white, colframe=blue!70!black,
  fonttitle=\bfseries, title={Tổng kết chương}, #1%
}


\begin{document}

\thispagestyle{empty}
\begin{tikzpicture}[remember picture, overlay]
  \fill[blue!80!black] (current page.south west) rectangle (current page.north east);
  \fill[blue!60!cyan, opacity=0.8]
    (current page.south west) -- (current page.north west) --
    ([xshift=-3cm]current page.north east) -- (current page.south east) -- cycle;
  \foreach \r/\o in {6/0.05, 4.5/0.08, 3/0.12} {
    \fill[white, opacity=\o] ([xshift=2cm, yshift=-2cm]current page.north east) circle (\r cm);
  }
  \draw[white, line width=2pt] ([yshift=-5cm]current page.north west) ++(2cm,0) -- ++(17cm,0);
  \draw[white, line width=2pt] ([yshift=-16cm]current page.north west) ++(2cm,0) -- ++(17cm,0);
  \node[anchor=west, white, font=\fontsize{28}{34}\bfseries\selectfont]
    at ([xshift=3cm, yshift=-7.5cm]current page.north west) {\coverbookline};
  \node[anchor=west, white, font=\fontsize{20}{26}\selectfont]
    at ([xshift=3cm, yshift=-9cm]current page.north west) {\covermaintitle};
  \node[anchor=west, white, font=\fontsize{16}{22}\selectfont]
    at ([xshift=3cm, yshift=-10.2cm]current page.north west) {\coversubtitle};
  \node[anchor=west, white, font=\Large\bfseries]
    at ([xshift=3cm, yshift=-13cm]current page.north west) {Biên soạn:};
  \node[anchor=west, yellow!80!white, font=\fontsize{24}{30}\bfseries\selectfont]
    at ([xshift=3cm, yshift=-14.5cm]current page.north west) {ĐẶNG TẤN QUÍ};
  \node[anchor=west, white!90, font=\large]
    at ([xshift=3cm, yshift=-18cm]current page.north west) {SĐT: 0377\,122\,210};
  \node[anchor=south, white!80, font=\small]
    at ([yshift=1.5cm]current page.south) {Tài liệu có bản quyền --- Cấm sao chép khi chưa được sự đồng ý của tác giả};
  \node[anchor=south east, white, font=\Large\bfseries]
    at ([xshift=-2cm, yshift=3cm]current page.south east) {\the\year};
\end{tikzpicture}

\newpage
\tableofcontents
\newpage

\section*{Lời nói đầu}
\addcontentsline{toc}{section}{Lời nói đầu}

Tài liệu thuộc \textbf{Nhóm 5 --- Xác suất và thống kê}, biên soạn theo cùng triết lý với giáo án Đại số tuyến tính: học để \emph{hiểu} và \emph{dùng}, không chỉ để ghi nhớ công thức.

\begin{enumerate}
  \item \textbf{Động cơ trước định nghĩa.} Mỗi phần lớn mở bằng ô \textsf{\color{teal!60!black} Động cơ học} --- câu hỏi thực tế hoặc bài toán mẫu cần khái niệm đó.
  \item \textbf{Trực giác trước công thức.} Ô \textsf{\color{purple!50!black} Ý nghĩa \& Trực giác} giải thích ``công thức đang nói gì'' (xác suất = tần suất dài hạn, kỳ vọng = trung bình có trọng số, v.v.).
  \item \textbf{Ví dụ có kiểm tra.} Ô \textsf{\color{violet!50!black} Ví dụ} nên tự làm trước khi đọc lời giải; phần thống kê ứng dụng cần gắn dữ liệu số cụ thể.
  \item \textbf{Tổng kết cuối mỗi chương.} Ô \textsf{\color{blue!70!black} Tổng kết chương} liệt kê kết quả phải nhớ và liên kết sang phần sau.
  \item \textbf{Phân tầng theo môn.} X1--X3 là nền; X4--X5 nâng cao lý thuyết; X6--X8 chuyên sâu (đa biến, Bayes, quá trình ngẫu nhiên).
\end{enumerate}

\noindent\textbf{Cách đọc:} Động cơ $\to$ Định nghĩa / Định lý $\to$ Ví dụ $\to$ Ý nghĩa $\to$ Tổng kết. Với môn thống kê, luôn hỏi: \emph{mẫu là gì, tham số nào chưa biết, giả thuyết $H_0$ là gì}.

\newpage


\section*{Tại sao học xác suất?}

\begin{dongco}[title={Động cơ học -- Ba tình huống quen thuộc}]
\textbf{Câu chuyện 1 -- Dự báo mưa.} App báo ``mưa 70\%'' chiều nay: không phải ``chắc chắn mưa'', cũng không phải ``đoán bừa''. Ta cần một \emph{số} trong $[0,1]$ để so sánh các kịch bản và ra quyết định (mang ô hay không). Đó là ý nghĩa của $P(A)$.

\medskip
\textbf{Câu chuyện 2 -- Xét nghiệm y tế.} Bệnh hiếm, test rất ``nhạy'' nhưng vẫn có người \textbf{dương tính giả}. Bác sĩ và bệnh nhân hay hiểu nhầm: ``dương tính'' $\not=$ ``chắc chắn mắc bệnh''. Công thức Bayes cho ta \emph{đảo ngược} xác suất: từ ``test dương'' suy ra ``thật sự mắc bệnh'' - cần cả tỷ lệ mắc bệnh trong dân số (tiên nghiệm).

\medskip
\textbf{Câu chuyện 3 -- Game và công bằng.} Tung đồng xu 10 lần, có người thắng 9 lần liên tiếp và bảo ``xu lệch''. Liệu xu có gian? Hay chỉ là \emph{biến động ngẫu nhiên}? Khi $n$ lớn, tần suất thắng ổn định quanh $1/2$ (luật số lớn) - nhưng với $n$ nhỏ, kết quả ``lệch'' vẫn bình thường.

\medskip
Ba câu chuyện trên cùng một ngôn ngữ: $\Omega$, sự kiện, $P(\cdot)$, có điều kiện, kỳ vọng. Phần đầu của tài liệu xây nền đó; các phần sau gán \emph{số} cho kết quả (biến ngẫu nhiên) và mô tả ``trung bình dài hạn''.
\end{dongco}

\newpage

\section{Sự kiện ngẫu nhiên và phép tính xác suất}

\begin{dongco}[title={Động cơ học -- Liệt kê hết mọi khả năng}]
Nói ``tung xu có thể sấp hoặc ngửa'' thì ai cũng hiểu. Nhưng ``rút 3 lá bài từ bộ 52'' - có bao nhiêu kết cục \emph{khác nhau}?

``Hẹn gặp trong khoảng 7h--8h, ai đến trước chờ 10 phút'' - kết cục là \emph{cặp thời gian} $(x,y)$, không phải một con số.

Trước khi tính $P(A)$, ta phải thống nhất: \textbf{không gian mẫu $\Omega$} là tập \emph{tất cả} kết cục có thể, không bỏ sót, không trùng lặp. Sự kiện $A$ chỉ là tập con của $\Omega$ - giống như lọc kết quả: ``mặt chẵn'' = $\{2,4,6\}$ trên xúc xắc.
\end{dongco}

\subsection{Phép thử và sự kiện}

\begin{dinhnghia}[title={Phép thử ngẫu nhiên và không gian mẫu}]
  \textbf{Phép thử ngẫu nhiên} (phép thử): Một phép thử có thể dẫn đến các \emph{kết cục} khác nhau dù được lặp lại trong điều kiện giống nhau.

  \textbf{Không gian mẫu} $\Omega$: tập tất cả kết cục có thể của 1 phép thử. Có thể liệt kê, biểu diễn bằng biểu thức, hoặc sơ đồ cây.
\end{dinhnghia}

\begin{vidu}[title={Không gian mẫu}]
\begin{itemize}
  \item Một đồng xu: $\Omega = \{H, T\}$ ($H$ = sấp, $T$ = ngửa), $|\Omega| = 2$.
  \item Ba đồng xu (có thứ tự): $\Omega = \{HHH, HHT, HTH, THH, HTT, THT, TTH, TTT\}$, $|\Omega| = 8$.
  \item Xúc xắc: $\Omega = \{1,2,3,4,5,6\}$, $|\Omega| = 6$.
  \item Đo độ dày đầu nối trong khoảng [10, 11] mm: $\Omega = \{x \in \mathbb{R} : 10 \le x \le 11\}$ (mm).
  \item Thời gian cuộc gọi ``dài hơn 5 phút'': $\Omega = \{x \mid x \ge 0\}$.
\end{itemize}
\end{vidu}

\begin{dinhnghia}[title={Sự kiện (biến cố)}]\textbf{Sự kiện} $A$ là tập con $A \subset \Omega$.
  \begin{itemize}
    \item \textbf{Sự kiện sơ cấp:} $\{\omega\}$ - chỉ một kết cục.
    \item \textbf{Sự kiện chắc chắn:} $\Omega$. \textbf{Không thể:} $\varnothing$.
    \item \textbf{Kéo theo:} $A \subset B$ nghĩa là $A$ xảy ra thì $B$ xảy ra; $A = B$ $\Leftrightarrow$ $A \subset B$ và $B \subset A$.
  \end{itemize}
\end{dinhnghia}

\begin{vidu}[title={Sự kiện trên xúc xắc}]
Gieo xúc xắc, $E_i$ = ``mặt $i$ chấm'', $i = 1,\ldots, 6$. $\Omega = \{E_1,\ldots,E_6\}$.
\begin{itemize}
  \item $A = \{E_2, E_4, E_6\}$: mặt chẵn.
  \item $B = \{E_4, E_5, E_6\}$: số chấm $> 3$.
  \item $E_2 \subset A$.
  \item Các $E_i$ \textbf{xung khắc từng đôi}; $\{E_1,\ldots,E_6\}$ là \textbf{hệ đầy đủ}.
\end{itemize}
\end{vidu}

\subsubsection{Phép toán trên sự kiện}

\begin{dinhnghia}
  \begin{itemize}
    \item \textbf{Hợp} $A \cup B$: ít nhất một trong hai xảy ra.
    \item \textbf{Giao} $A \cap B$ (viết $AB$): cả hai cùng xảy ra.
    \item \textbf{Hiệu} $A \setminus B = A \cap \bar{B}$: $A$ xảy ra, $B$ không.
    \item \textbf{Đối / phần bù} $\bar{A} = \Omega \setminus A$.
    \item \textbf{Xung khắc:} $A \cap B = \varnothing$.
    \item \textbf{Hệ đầy đủ} $\{B_1,\ldots,B_n\}$: đôi một xung khắc và $\bigcup B_i = \Omega$.
    \item \textbf{Hệ xung khắc từng đôi} $\{A_1,\ldots,A_n\}$: $A_i \cap A_j = \varnothing$ với $i \neq j$.
  \end{itemize}
\end{dinhnghia}

\begin{vidu}[title={Hợp và giao - mạch điện}]
Ba bóng $i = 1,2,3$, $A_i$ = ``bóng $i$ cháy'', $A$ = ``mạch mất điện''.
\begin{itemize}
  \item \textbf{Mắc nối tiếp:} mất điện khi \emph{ít nhất một} bóng cháy $\Rightarrow$ $A = A_1 \cup A_2 \cup A_3$.
  \item \textbf{Mắc song song:} mất điện khi \emph{cả ba} cháy $\Rightarrow$ $A = A_1 \cap A_2 \cap A_3$.
\end{itemize}
\end{vidu}

\begin{tinhchat}[title={Đại số tập hợp}]
\begin{itemize}
  \item Giao hoán: $A \cup B = B \cup A$, $A \cap B = B \cap A$.
  \item Kết hợp: $A \cup (B \cup C) = (A \cup B) \cup C$, $A \cap (B \cap C) = (A \cap B) \cap C$.
  \item Phân phối: $A \cup (B \cap C) = (A \cup B) \cap (A \cup C)$, $A \cap (B \cup C) = (A \cap B) \cup (A \cap C)$.
  \item De Morgan: $\overline{A \cup B} = \bar{A} \cap \bar{B}$, $\overline{A \cap B} = \bar{A} \cup \bar{B}$.
  \item $A \cup A = A$, $A \cap A = A$.
  \item $A \cup \varnothing = A$, $A \cap \varnothing = \varnothing$.
  \item $A \cup \Omega = \Omega$, $A \cap \Omega = A$.
  \item $A \cup \bar{A} = \Omega$, $A \cap \bar{A} = \varnothing$.
  \item $\bar{\Omega} = \varnothing$, $\bar{\varnothing} = \Omega$.
  \item $\bar{\bar{A}} = A$.
\end{itemize}
\end{tinhchat}

\begin{vidu}[title={Hệ đầy đủ - ba phân xưởng}]
Chọn ngẫu nhiên một sản phẩm; $C_i$ = ``do phân xưởng $i$ sản xuất'', $i = 1,2,3$. Hệ $\{C_1, C_2, C_3\}$ đầy đủ (mỗi sản phẩm thuộc đúng một phân xưởng).
\end{vidu}

\subsection{Phương pháp đếm}

\begin{dongco}[title={Động cơ học -- Mật khẩu và vé số}]
\textbf{Mật khẩu 4 chữ số} (mỗi vị trí chọn $0$--$9$): có $10^4 = 10\,000$ mật khẩu - quy tắc \textbf{nhân} (giai đoạn 1 chọn hàng nghìn, giai đoạn 2 chọn hàng trăm, \ldots).

\textbf{Chọn 3 người làm đội trưởng trong lớp 40 SV} (không trùng vai): $40 \times 39 \times 38$ cách - có \textbf{thứ tự}, không lặp ($A_{40}^3$). Nếu chỉ cần ``bộ ba người'' không phân vai: chia $3!$ vì đếm trùng - ra $C_{40}^3$.

Hai bài toán tưởng giống nhau nhưng đếm khác nhau vì \emph{``có quan tâm thứ tự hay không?''}. Phần này cho bộ công cụ đếm để sau đó viết $P(A) = |A|/|\Omega|$ mà không sót, không đếm đôi.
\end{dongco}

\begin{phuongphap}[title={Quy tắc 1 và 2}]
\textbf{Cộng:} Nếu một công việc có $k$ phương án khác nhau để thực hiện, phương án $i$ có $n_i$ cách thực hiện xong công việc và không có một cách thực hiện nào ở phương án này lại trùng với một cách thực hiện ở phương án khác. $\Rightarrow$ tổng $n_1 + n_2 + \cdots + n_k$ cách.

\textbf{Nhân:} Giả sử một công việc nào đó được chia thành $k$ giai đoạn. Có $n_i$ cách thực hiện giai đoạn thứ $i$ thì $\Rightarrow$ $n_1 \times n_2 \times \cdots \times n_k$ cách.
\end{phuongphap}

\begin{vidu}[title={Đi từ A đến C}]
Giả sử để đi từ A đến C có thể đi qua B, trong đó có 2 đường khác nhau đi trực tiếp từ A đến C, có 3 đường khác nhau để đi từ A đến B và có 2 đường khác nhau để đi từ B đến C. Hỏi có bao nhiêu cách đi từ A đến C?

Tổng: $2 + 6 = 8$ cách.
\end{vidu}

\begin{congthuc}[title={Chỉnh hợp, chỉnh hợp lặp, hoán vị, tổ hợp}]
\begin{center}
\renewcommand{\arraystretch}{1.15}
\begin{tabular}{lll}
\toprule
Khái niệm & Số cách & Ghi chú \\
\midrule
Chỉnh hợp chập $k$ & $A_n^k = \dfrac{n!}{(n-k)!}$ & Có thứ tự, không lặp \\
Chỉnh hợp lặp chập $k$ & $\bar{A}_n^k = n^k$ & Có thứ tự, cho lặp \\
Hoán vị & $P_n = n!$ & $A_n^n$ \\
Hoán vị có lặp & $\dfrac{n!}{n_1! n_2! \cdots n_k!}$ & $n = n_1 + \cdots + n_k$ \\
Tổ hợp chập $k$ & $C_n^k = \dfrac{n!}{k!(n-k)!}$ & Không thứ tự \\
\bottomrule
\end{tabular}
\end{center}
$C_n^k = A_n^k / k!$.
\end{congthuc}

\begin{vidu}[title={Ví dụ đếm}]
\begin{itemize}
  \item Một bảng mạch in có tám vị trí khác nhau, mỗi vị trí có thể đặt được một thành phần của mạch. Nếu bốn thành phần khác nhau được đặt trên bảng mạch thì có thể có bao nhiêu kiểu dáng khác nhau?: $A_8^4 = 8 \cdot 7 \cdot 6 \cdot 5 = 1680$ kiểu dáng.
  \item Từ năm chữ số $1, 2, 3, 4, 5$ lập được bao nhiêu số có ba chữ số?: $\bar{A}_5^3 = 125$ (lặp); nếu khác nhau: $A_5^3 = 60$.
  \item Một sản phẩm được dán nhãn bằng cách in năm dòng dày, ba dòng trung bình và hai dòng mỏng. Nếu mỗi thứ tự của mười dòng đại diện cho một nhãn khác nhau, thì có thể tạo ra bao nhiêu nhãn khác nhau bằng cách sử dụng phương pháp này?: $\dfrac{10!}{5!\,3!\,2!} = 2520$.
  \item Một lô hàng chứa 60 sản phẩm, trong đó có 5 sản phẩm bị lỗi. Lấy ngẫu nhiên từ lô hàng ra 6 sản phẩm. Có bao nhiêu cách khác nhau để lấy được đúng 2 sản phẩm bị lỗi?: $C_5^2 C_{55}^4$ cách.
  \item Có 6 người khách được xếp vào 6 ghế quanh một bàn tròn có 6 chỗ.
  \begin{itemize}
    \item Nếu các ghế ngồi đã được đánh số từ 1 đến 6 thì có bao nhiêu cách sắp xếp?: $P_6 = 6! = 720$ cách xếp.
    \item Nếu các ghế ngồi không được đánh số thứ tự thì sẽ có bao nhiêu cách? $P_5 = 5! = 120$ cách xếp.
  \end{itemize}
\end{itemize}
\end{vidu}

\subsection{Xác suất của một sự kiện}

\begin{dongco}
Ta đã biết liệt kê $\Omega$ và sự kiện $A$. Bước tiếp: gán cho mỗi sự kiện một \textbf{con số} $P(A) \in [0,1]$ để so sánh ``dễ xảy ra'' hơn hay kém hơn.

Tung đồng xu \emph{công bằng}: mỗi mặt $1/2$. Xúc xắc lỗi (mặt 6 nặng hơn): các mặt \emph{không} còn đồng khả năng - khi đó không dùng $|A|/|\Omega|$ đơn giản, mà cần \emph{mô hình} hoặc \emph{tần suất} từ nhiều lần lặp. Ba quan điểm dưới (cổ điển, hình học, thống kê) là ba cách ``chọn'' con số $P(A)$ phù hợp với bối cảnh.
\end{dongco}

\begin{dinhnghia}[title={Khái niệm xác suất}]
Tiến hành phép thử, $A \subset \Omega$. \textbf{Xác suất} $P(A) \in [0,1]$ đo mức ``chắc chắn'' $A$ xảy ra. $P(\Omega) = 1$, $P(\varnothing) = 0$.
\end{dinhnghia}

\subsubsection{Quan điểm cổ điển}
\begin{dinhnghia}[title=Kết cục đồng khả năng]
Các kết cục của một phép thử được gọi là đồng khả năng nếu khả năng xảy ra của chúng là như nhau.

Chẳng hạn, trong một phép thử có n kết cục có thể xảy ra, các kết cục này là đồng khả năng nếu khả năng xảy ra của mỗi kết cục đều là $1/n$.
\end{dinhnghia}

\begin{dinhnghia}
  Giả sử trong một phép thử có n kết cục đồng khả năng, trong đó, có m kết cục thuận lợi cho sự xuất hiện của sự kiện A. Khi đó:
  \[
    P(A) = \frac{\text{số kết cục thuận lợi cho } A}{\text{tổng số kết cục đồng khả năng}}
    = \frac{|A|}{|\Omega|} = \frac{m}{n}.
  \]
\end{dinhnghia}

\begin{tinhchat}[title={Tính chất của xác suất}]
\begin{itemize}
  \item $0 \le P(A) \le 1$.
  \item $P(\Omega) = 1$.
  \item $P(\varnothing) = 0$.
  \item $A \subset B \Rightarrow P(A) \le P(B)$.
  \item $P(\bar{A}) = 1 - P(A)$.
\end{itemize}
\end{tinhchat}

\begin{vidu}[title={Gọi điện quên 2 số cuối}]
Nhớ hai số cuối \emph{khác nhau}, chọn ngẫu nhiên: $|\Omega| = A_{10}^2 = 90$, chỉ $1$ cách đúng $\Rightarrow P(A) = 1/90$.
\end{vidu}

\begin{vidu}[title={Rút 2 lá bài}]
52 lá, rút 2 lá không hoàn lại: $|\Omega| = C_{52}^2 = 1326$.
\begin{itemize}
  \item Cả hai J: $P(A) = C_4^2 / 1326 = 6/1326 = 1/221$.
  \item Một J một K: $P(B) = C_4^1 C_4^1 / 1326 = 16/1326 = 8/663$.
\end{itemize}
\end{vidu}

\subsubsection{Quan điểm hình học}

\begin{dongco}
  Định nghĩa xác suất theo quan điểm cổ điển có ưu điểm là dễ vận dụng, tuy nhiên, định nghĩa này chỉ áp dụng được với các phép thử ngẫu nhiên có hữu hạn kết cục đồng khả năng. Trong trường hợp có vô hạn kết cục đồng khả năng ta sử dụng định nghĩa xác suất theo quan điểm hình học.
\end{dongco}

\begin{dinhnghia}
Giả sử một phép thử có vô hạn kết cục đồng khả năng và các kết cục này được biểu thị bởi một miền hình học $G$ (có độ đo hữu hạn và khác $0$), còn các kết cục thuận lợi cho sự kiện $A$ được biểu thị bởi miền con $H$ của $G$, với $H \subset G$. Khi đó:
  \[
    P(A) = \frac{\mu(H)}{\mu(G)}
  \]
  ($\mu$ = độ dài / diện tích / thể tích).
\end{dinhnghia}

\begin{vidu}[title={Hẹn gặp 7h-8h, chờ 10 phút}]
Hai người bạn hẹn gặp nhau ở một địa điểm trong khoảng thời gian từ 7h00 đến 8h00. Mỗi người có thể đến điểm hẹn một cách ngẫu nhiên tại một thời điểm trong khoảng thời gian nói trên và họ quy ước rằng ai đến mà không thấy người kia thì chỉ đợi trong vòng 10 phút. Tính xác suất để hai người gặp nhau.

$(x,y)$ = thời điểm hai người đến (phút, $0 \le x,y \le 60$). $G$ = hình vuông cạnh $60$. Gặp nhau khi $|x - y| \le 10$ (dải quanh đường chéo).
\[
  P(E) = \dfrac{60^2 - 50^2}{60^2} = \dfrac{11}{36} \approx 0{,}306.
\]
\end{vidu}

\subsubsection{Quan điểm thống kê}

\begin{dongco}
Cổ điển và hình học đều cần \textbf{đồng khả năng} - trong thực tế thường khó kiểm tra; khi đó dùng quan điểm thống kê hoặc mô hình xác suất chung.
\end{dongco}

\begin{dinhnghia}
Giả sử trong một điều kiện nào đó ta có thể lặp lại $n$ lần một phép thử và thấy có $x$ lần xuất hiện sự kiện $A$. Khi đó, $x$ được gọi là tần số và tỷ số $x/n$ gọi là tần suất xuất hiện sự kiện $A$, ký hiệu là $f_n(A)$. Như vậy:
\[
f_n(A) = \dfrac{x}{n}
\]
Khi số phép thử tăng lên vô hạn, tần suất xuất hiện sự kiện $A$ tiến dần đến một số xác định $p$ nào đó. Số $p$ đó được gọi là xác suất của sự kiện $A$ (theo quan điểm thống kê).
\[
P(A) = \lim_{n \to \infty} f_n(A) = p
\]
\end{dinhnghia}

\begin{vidu}[title={Tung đồng xu}]
Để xác định tần suất xuất hiện mặt sấp khi tung một đồng xu nhiều lần, người ta ghi lại kết quả sau:
\begin{center}
\begin{tabular}{|c|c|c|c|}
Người thí nghiệm & Số lần tung (n) & Số lần xuất hiện mặt sấp (x) & Tần suất (fn) \\
\hline
Buffon & 4040 & 2048 & 0,5069 \\
Pearson & 12000 & 6019 & 0,5016 \\
Pearson & 24000 & 12012 & 0,5005 \\
\end{tabular}
\end{center}
Ta thấy rằng khi số lần tung đồng xu càng lớn, tần suất xuất hiện mặt sấp càng gần với xác suất xuất hiện mặt sấp là $0.5$.
\end{vidu}

\begin{luuy}
  \begin{itemize}
    \item Định nghĩa xác suất theo quan điểm thống kê khắc phục được một nhược điểm của định nghĩa xác suất theo quan điểm cổ điển là không dùng đến khái niệm đồng khả năng.
    \item Định nghĩa xác suất theo quan điểm thống kê không giúp ta tính được chính xác xác suất của một sự kiện mà chỉ tìm được giá trị gần đúng; đồng thời số phép thử phải đủ lớn và chỉ dùng được cho các phép thử có thể lặp lại nhiều lần một cách độc lập trong các điều kiện giống nhau.
  \end{itemize}
\end{luuy}

\subsection{Công thức cộng và nhân xác suất}

\begin{dongco}
\textbf{``Hoặc'' và ``và'' trên sự kiện.} Lớp 100 SV: 40 giỏi ngoại ngữ, 30 giỏi toán, 20 giỏi \emph{cả hai}. Chọn ngẫu nhiên 1 người - xác suất giỏi \emph{ít nhất một môn} là bao nhiêu?

Nếu cộng thô $40\% + 30\% = 70\%$ thì \textbf{sai}, vì 20\% bị tính hai lần. Công thức cộng $P(N \cup T) = P(N) + P(T) - P(N \cap T)$ chính là cách \emph{trừ phần giao bị đếm trùng}.

\medskip
\textbf{``Biết rằng'' thay đổi xác suất.} Rút bài: xác suất rút được J là $4/52$. Nhưng nếu \emph{biết} lá vừa rút là đen, chỉ còn 26 lá đen trong đó có 2 J $\Rightarrow$ $P(J \mid \text{đen}) = 2/26$. Công thức nhân và xác suất có điều kiện mô tả chính xác kiểu suy luận ``cập nhật tin'' này.
\end{dongco}

\begin{dinhly}[title={Công thức cộng}]
\[
  P(A \cup B) = P(A) + P(B) - P(A \cap B).
\]
Nếu $A$ và $B$ xung khắc thì $A \cap B = \varnothing$: $P(A \cup B) = P(A) + P(B)$.

\[
  P(\bigcup_{i=1}^{n} A_i) = \sum_{i=1}^{n} P(A_i) - \sum_{1 \le i < j \le n} P(A_iA_j) + \cdots + (-1)^{n-1} P(A_1A_2 \cdots A_n).
\]

Nếu $A_1, \ldots, A_n$ đôi một xung khắc thì: $P(\bigcup A_i) = \sum P(A_i)$.
\end{dinhly}

\begin{vidu}[title={Lớp 100 sinh viên}]
Một lớp có 100 sinh viên, trong đó có 40 sinh viên giỏi ngoại ngữ, 30 sinh viên giỏi toán, 20 sinh viên vừa giỏi ngoại ngữ, vừa giỏi toán. Chọn ngẫu nhiên một sinh viên trong lớp. Tìm xác suất để sinh viên đó giỏi ít nhất một trong hai môn trên:

Gọi $N$ là sự kiện “sinh viên đó giỏi ngoại ngữ”; $T$ là sự kiện “sinh viên đó giỏi toán”.
\[
  P(N \cup T) = P(N) + P(T) - P(N \cap T) = \frac{40}{100} + \frac{30}{100} - \frac{20}{100} = \frac{50}{100} = 0{,}5
\]
\end{vidu}
\begin{dinhnghia}[title={Xác suất có điều kiện}]
Giả sử trong một phép thử có sự kiện B với P(B) > 0. Khi đó, xác suất có điều kiện của sự kiện A nào đó,
biết rằng sự kiện B đã xảy ra là:
\[
  \displaystyle P(A \mid B) = \frac{P(A \cap B)}{P(B)}
\]

Tương tự: $P(B \mid A) = \frac{P(A \cap B)}{P(A)}$.
\end{dinhnghia}
\begin{vidu}[title={Rút bài -- biết quân đen}]
Từ một bộ bài tú lơ khơ gồm 52 quân bài đã được trộn kỹ rút ngẫu nhiên ra một quân bài. Biết rằng quân bài được rút ra là quân màu đen, tính xác suất đó là quân J.

Gọi $A$ là sự kiện “rút được quân J” và $B$ là sự kiện “rút được quân màu đen”.

$P(A \mid B) = \dfrac{2}{26} = \dfrac{1}{13}$.

Hoặc $P(A \mid B) = \dfrac{P(AB)}{P(B)} = \dfrac{2/52}{26/52} = \dfrac{1}{13}$.
\end{vidu}

\begin{dinhnghia}[title={Độc lập}]
$A$, $B$ \textbf{độc lập} nếu sự kiện này xuất hiện hay không xuất hiện không làm ảnh hưởng tới khả năng xuất hiện của sự kiện kia, nghĩa là $P(A \mid B) = P(A \mid \bar{B}) = P(A)$ và $P(B \mid A) = P(B \mid \bar{A}) = P(B)$.

Các sự kiện $A_1, A_2, \ldots, A_n$ được gọi là \textbf{độc lập từng đôi} với nhau nếu mỗi cặp hai trong $n$ sự kiện đó độc lập với nhau.

Các sự kiện $A_1, A_2, \ldots, A_n$ được gọi là \textbf{độc lập tổng thể} nếu mỗi sự kiện trong chúng độc lập với tích của một số $s$ bất kỳ sự kiện trong các sự kiện còn lại, $1 \le s \le n - 1$.
\end{dinhnghia}

\begin{phanvidu}[title={Không độc lập tổng thể}]
Gieo đồng thời hai đồng xu cân đối đồng chất lên mặt phẳng cứng. Gọi $A_1$ là sự kiện “ít nhất một đồng xu
xuất hiện mặt sấp”, $A_2$ là sự kiện “ít nhất một đồng xu xuất hiện mặt ngửa”, $A_3$ là sự kiện “có ba đồng xu
xuất hiện mặt ngửa”. $A_1, A_2, A_3$ có độc lập tổng thể không?

$P(A_1 A_2 A_3) = P(A_1)P(A_2)P(A_3) = 0$ nhưng $P(A_1 A_2) = 1/2 \ne 9/16 = P(A_1)P(A_2)$

Đó đó \textbf{không} độc lập tổng thể.
\end{phanvidu}

\begin{dinhly}[title={Công thức nhân}]
$P(AB) = P(A)\,P(B \mid A) = P(B)\,P(A \mid B)$.

Nếu A và B độc lập thì $P(AB) = P(A)P(B)$.

$P(A_1 A_2 \cdots A_n) = P(A_1)\,P(A_2 \mid A_1)\cdots P(A_n \mid A_1 \cdots A_{n-1})$.

Nếu $A_1, A_2, \ldots, A_n$ độc lập tổng thể thì $P(A_1 A_2 \cdots A_n) = P(A_1)\,P(A_2)\cdots P(A_n)$.
\end{dinhly}

\begin{vidu}[title={Ba xạ thủ}]
Ba xạ thủ cùng bắn súng vào một bia độc lập nhau. Xác suất bắn trúng bia của xạ thủ thứ nhất, thứ hai, thứ ba tương ứng là 0,7, 0,8 và 0,9. Tính xác suất để:
\begin{itemize}
  \item Có duy nhất một xạ thủ bắn trúng bia;
  \item Có đúng hai xạ thủ bắn trúng bia;
  \item Có ít nhất một xạ thủ bắn trúng bia;
  \item Xạ thủ thứ nhất bắn trúng bia biết rằng có hai xạ thủ bắn trúng bia
\end{itemize}

Giải:

\begin{itemize}
  \item Có duy nhất một xạ thủ bắn trúng bia:
    $P(A) = P(A_1 \bar{A_2} \bar{A_3}) + P(\bar{A_1} A_2 \bar{A_3}) + P(\bar{A_1} \bar{A_2} A_3) = 0{,}7 \cdot 0{,}2 \cdot 0{,}1 + 0{,}3 \cdot 0{,}8 \cdot 0{,}1 + 0{,}3 \cdot 0{,}2 \cdot 0{,}9 = 0{,}096$
  \item Có đúng hai xạ thủ bắn trúng bia:
    $P(B) = P(A_1 A_2 \bar{A_3}) + P(A_1 \bar{A_2} A_3) + P(\bar{A_1} A_2 A_3) = 0{,}7 \cdot 0{,}8 \cdot 0{,}1 + 0{,}7 \cdot 0{,}2 \cdot 0{,}9 + 0{,}3 \cdot 0{,}8 \cdot 0{,}9 = 0{,}398$
  \item Có ít nhất một xạ thủ bắn trúng bia:
    $P(C) = P(A_1 \cup A_2 \cup A_3) = P(A_1) + P(A_2) + P(A_3) - P(A_1 A_2) - P(A_1 A_3) - P(A_2 A_3) + P(A_1 A_2 A_3) = 0{,}7 + 0{,}8 + 0{,}9 - 0{,}7 \cdot 0{,}8 - 0{,}7 \cdot 0{,}9 - 0{,}8 \cdot 0{,}9 + 0{,}7 \cdot 0{,}8 \cdot 0{,}9 = 0{,}994$

    Hoặc: $P(C) = 1 - P(\bar{A_1} \bar{A_2} \bar{A_3}) = 1 - 0{,}3 \cdot 0{,}2 \cdot 0{,}1 = 0{,}994$

    Hoặc: $P(C) = P(A) + P(B) + P(A_1 A_2 A_3) = 0{,}096 + 0{,}398 + 0{,}7 \cdot 0{,}8 \cdot 0{,}9 = 0{,}994$
  \item Xạ thủ thứ nhất bắn trúng bia biết rằng có hai xạ thủ bắn trúng bia:
    $P(A_1 \mid B) = \dfrac{P(A_1 B)}{P(B)} = \dfrac{P(A_1 A_2 \bar{A_3}) + P(A_1 \bar{A_2} A_3)}{P(B)} = \dfrac{0{,}7 \cdot 0{,}8 \cdot 0{,}1 + 0{,}7 \cdot 0{,}2 \cdot 0{,}9}{0{,}398} = 0{,}46$
\end{itemize}
\end{vidu}

\subsubsection{Dãy Bernoulli và công thức Bernoulli}

\begin{dinhnghia}
  \textbf{Dãy $n$ phép thử Bernoulli:} $n$ lần độc lập, mỗi lần chỉ ``thành công'' / ``thất bại'', xác suất thành công $p$ không đổi.
\end{dinhnghia}

\begin{dinhly}[title={Công thức Bernoulli}]
Xác suất đúng $k$ lần thành công trong $n$ lần:
\[
  P_n(k) = C_n^k\, p^k (1-p)^{n-k}, \qquad k = 0, 1, \ldots, n.
\]
\end{dinhly}

\begin{vidu}[title={Bắn cung}]
$p = 0{,}8$, $n = 4$, trúng đúng $3$ lần: $P_4(3) = C_4^3 (0{,}8)^3 (0{,}2) = 0{,}4096$.
\end{vidu}

\subsection{Xác suất toàn phần và công thức Bayes}

\begin{dongco}[title={Động cơ học -- Sản phẩm đến từ phân xưởng nào?}]
Lô hàng gồm hàng từ ba phân xưởng với tỷ lệ khác nhau; mỗi phân xưởng có tỷ lệ sản phẩm loại 1 khác nhau. Bạn \emph{nhìn thấy} một sản phẩm loại 1 - hỏi: nó \textbf{khả năng cao} đến từ phân xưởng nào?

Đây là bài toán \emph{ngược}: đã biết kết quả quan sát ($A$), muốn suy đoán \emph{nguyên nhân} ($B_k$).

Với xét nghiệm y tế: ``dương tính'' không đồng nghĩa ``chắc chắn bệnh'' nếu bệnh hiếm.
\end{dongco}

\begin{dinhly}[title={Xác suất toàn phần}]
$\{B_1, \ldots, B_n\}$ hệ đầy đủ, $P(B_i) > 0$. Với mọi sự kiện $A$:
\[
  P(A) = \sum_{i=1}^{n} P(AB_i) = \sum_{i=1}^{n} P(B_i)\,P(A \mid B_i).
\]
\end{dinhly}

\begin{vidu}[title={Ba phân xưởng}]
Lô hàng: 20\% PX1, 50\% PX2, 30\% PX3. Tỷ lệ sản phẩm loại 1: $0{,}7$, $0{,}8$, $0{,}6$. $B_i$ = ``lấy từ PX$i$'', $A$ = ``loại 1'':
\[
  P(A) = 0{,}2 \cdot 0{,}7 + 0{,}5 \cdot 0{,}8 + 0{,}3 \cdot 0{,}6 = 0{,}72.
\]
\end{vidu}

\begin{dinhly}[title={Công thức Bayes}]
\[
  P(B_k \mid A) = \frac{P(B_k)\,P(A \mid B_k)}{P(A)} = \frac{P(B_k)\,P(A \mid B_k)}{\displaystyle\sum_{i=1}^{n} P(B_i)\,P(A \mid B_i)}.
\]
\end{dinhly}

\begin{vidu}[title={Sản phẩm loại 1 do PX nào?}]
Với $P(A) = 0{,}72$ như trên:
\[
  P(B_1 \mid A) = \frac{0{,}2 \cdot 0{,}7}{0{,}72} \approx 0{,}194,\quad
  P(B_2 \mid A) \approx \frac{0{,}4}{0{,}72} \approx 0{,}556,\quad
  P(B_3 \mid A) = 0{,}25.
\]
Khả năng cao nhất: \textbf{phân xưởng II} ($P(B_2 \mid A)$ lớn nhất).
\end{vidu}

\begin{ungdung}[title={Bayes trong chẩn đoán}]
Bệnh hiếm $P(B) = 0{,}01$, test nhạy $P(+ \mid B) = 0{,}99$, đặc hiệu $P(- \mid \bar{B}) = 0{,}95$. Người \textbf{dương tính} thật sự mắc bệnh:
\[
  P(B \mid +) = \frac{0{,}99 \cdot 0{,}01}{0{,}99 \cdot 0{,}01 + 0{,}05 \cdot 0{,}99} \approx 0{,}17.
\]
Dương tính không đồng nghĩa chắc chắn mắc bệnh.
\end{ungdung}

\begin{luuy}
Muốn dùng toàn phần / Bayes phải chọn \textbf{hệ đầy đủ thích hợp}; sự kiện cần tính $P(B_k \mid A)$ phải là một $B_k$ trong hệ đó.
\end{luuy}

\begin{tongket}[title={Tổng kết}]
\begin{enumerate}
  \item $\Omega$, sự kiện, $\cup,\cap,\bar{\ },\setminus$, hệ đầy đủ, De Morgan.
  \item Đếm: $n_1 + \cdots + n_k$; $n_1 \cdots n_k$; $A_n^k$, $C_n^k$, $n!$, hoán vị lặp.
  \item $P(A)$: cổ điển $|A|/|\Omega|$; hình học $\mu(H)/\mu(G)$; tần suất $f_n(A)$.
  \item Cộng, có điều kiện, nhân, độc lập; Bernoulli $C_n^k p^k(1-p)^{n-k}$.
  \item Toàn phần $\sum P(B_i)P(A \mid B_i)$; Bayes; tiên nghiệm / hậu nghiệm.
\end{enumerate}

\end{tongket}

%% ================================================================
%%  PHẦN II--IV
%% ================================================================

\section{Biến ngẫu nhiên và luật phân phối}

\begin{dongco}
Trong thực tế ta thường gán số: điểm thi, chiều cao, số lỗi trên dây chuyền, thời gian chờ cuộc gọi.

Một hàm $X(\omega)$ từ $\Omega$ vào $\mathbb{R}$ gọi là \textbf{biến ngẫu nhiên}. Thay vì hỏi ``kết cục nào'', ta hỏi ``$X$ bằng bao nhiêu, với xác suất bao nhiêu?'' - toàn bộ thông tin gom vào \textbf{phân phối} ($p_X$, $F_X$, hoặc $f_X$).

\textbf{Điểm thi trắc nghiệm.} 100 câu, mỗi câu đúng/sai - số câu đúng $X$ chỉ nhận $0,1,\ldots,100$ (rời rạc). \textbf{Chiều cao} sinh viên trong lớp: có thể là $162{,}3$ cm, $175{,}8$ cm\ldots (liên tục trên khoảng).

Cùng là ``đại lượng ngẫu nhiên'', nhưng cách mô tả khác nhau: bảng $P(X=k)$ cho rời rạc, hàm mật độ cho liên tục. Phần này xây bộ ngôn ngữ đó trước khi đi vào các mô hình cụ thể (nhị thức, Poisson, chuẩn\ldots).
\end{dongco}

\subsection{Khái niệm và phân loại biến ngẫu nhiên}

\begin{dinhnghia}
  Cho phép thử với không gian mẫu $\Omega$. \textbf{Biến ngẫu nhiên (BNN)} là một hàm số của các kết cục $X: \Omega \to \mathbb{R}$, $\omega \mapsto X(\omega)$.
  \begin{itemize}
    \item $S_X = \{X(\omega) : \omega \in \Omega\}$: tập giá trị.
    \item \textbf{Rời rạc:} $S_X$ hữu hạn hoặc vô hạn đếm được.
    \item \textbf{Liên tục:} $S_X$ vô hạn không đếm được
    \item \textbf{Hỗn hợp:} $S_X$ vừa rời rạc vừa liên tục (lược bỏ).
  \end{itemize}
  Sự kiện $(X = x_i)$ tạo hệ đầy đủ khi $X$ chỉ nhận $x_1,\ldots,x_n$ trên một phép thử.
\end{dinhnghia}

\begin{vidu}[title={Số mặt ngửa khi tung xu 2 lần}]
$\Omega = \{SS, SN, NS, NN\}$. $X$ = số lần ngửa $\Rightarrow$ $S_X = \{0,1,2\}$:
$SS \mapsto 0$; $SN, NS \mapsto 1$; $NN \mapsto 2$.
\end{vidu}

\subsubsection{Hàm của một biến ngẫu nhiên}

\begin{dinhnghia}
Cho $X$ là một biến ngẫu nhiên và $g(x)$ là một hàm số thực. Đặt $Y = g(X)$. Biến ngẫu nhiên $g(X)$ như vậy được gọi là hàm của biến ngẫu nhiên $X$.

Nếu $X$ là biến ngẫu nhiên rời rạc thì $g(X)$ cũng là biến ngẫu nhiên rời rạc.

Nếu $X$ là biến ngẫu nhiên liên tục và $g(x)$ là một hàm liên tục thì $g(X)$ là biến ngẫu nhiên liên tục.
\end{dinhnghia}

\begin{vidu}[title={Hàm của biến ngẫu nhiên}]
Một công ty điện thoại tính cước phí cho một bản fax như sau: 10 nghìn đồng cho trang thứ nhất, 9 nghìn đồng cho trang thứ hai, ..., 6 nghìn đồng cho trang thứ năm. Những bản fax từ 6 đến 10 trang có phí là 50 nghìn đồng. Công ty điện thoại này không nhận những bản fax dài quá 10 trang. Gọi $X$ là số trang trong một bản fax và $Y$ là chi phí phải trả cho một bản fax. Khi đó, $X$ là biến ngẫu nhiên rời rạc và $S_X = \{1, 2, ..., 10\}$ và $Y$ là một hàm của $X$ được xác định bởi
  $Y = g(X) = \begin{cases} 10.5X - 0.5X^2, & 1 \le X \le 5 \\ 50, & 6 \le X \le 10 \end{cases}$
\end{vidu}

\subsection{Phân phối xác suất của biến ngẫu nhiên}

\begin{dongco}
Cho $P(A)$ với $A$ là tập con $\Omega$. Với BNN, ta quan tâm sự kiện dạng $(X = k)$, $(X \le x)$, $(a < X \le b)$.

\textbf{Hàm xác suất} $p_X(k)$ (rời rạc) hoặc \textbf{hàm phân phối} $F_X(x) = P(X \le x)$ (mọi loại) là ``bản đồ'' đầy đủ: biết $F_X$ là biết mọi xác suất dạng khoảng. Với liên tục, $f_X$ là đạo hàm của $F_X$ - giống ``mật độ'' trên trục số: vùng cao thì $X$ hay rơi vào đó hơn.
\end{dongco}

\subsubsection{Hàm xác suất (rời rạc)}

\begin{dinhnghia}
  BNN rời rạc $X$ nhận $x_1, x_2, \ldots$. \textbf{Hàm xác suất} $p_X(x)$:
  \begin{itemize}
    \item $p_X(x_i) \ge 0$;
    \item $\sum_i p_X(x_i) = 1$;
    \item $p_X(x_i) = P(X = x_i)$.
  \end{itemize}
\end{dinhnghia}

\begin{vidu}[title={Số bít lỗi trong 4 bít}]
Một bít được truyền qua đường truyền kỹ thuật số có thể bị lỗi. Xác suất để một bít được truyền đi bị lỗi là $0.1$. Giả sử rằng các lần truyền là độc lập nhau. Gọi $X$ là số bít bị lỗi trong bốn bít được truyền đi. Khi đó, $X$ là biến ngẫu nhiên rời rạc và $S_X = \{0, 1, 2, 3, 4\}$. Áp dụng công thức Bernoulli
\[
\begin{array}{c|ccccc}
k & 0 & 1 & 2 & 3 & 4 \\
\hline
P(X=k) & 0.6561 & 0.2916 & 0.0486 & 0.0036 & 0.0001
\end{array}
\]
\end{vidu}

\subsubsection{Hàm phân phối}

\begin{dinhnghia}
Hàm phân phối xác suất của biến ngẫu nhiên $X$, ký hiệu là $F_X(x)$, được định nghĩa như sau: $F_X(x) = P(X \le x)$, $x \in \mathbb{R}$.
\end{dinhnghia}

\begin{tinhchat}[title={Tính chất của hàm phân phối}]
\begin{itemize}
  \item $0 \le F_X(x) \le 1$ với mọi $x \in \mathbb{R}$.
  \item $F_X(x)$ không giảm và liên tục trái và liên tục phải.
  \item $P(a < X \le b) = P(a \le X < b) = P(a \le X \le b) = P(a < X < b) = F_X(b) - F_X(a)$.
  \item $\lim_{x \to -\infty} F_X(x) = 0$, $\lim_{x \to +\infty} F_X(x) = 1$.
\end{itemize}
\end{tinhchat}

\begin{vidu}[title={Hàm phân phối của biến ngẫu nhiên}]
  Cho $X$ là biến ngẫu nhiên rời rạc với bảng phân phối xác suất:
  \begin{center}
    \begin{tabular}{c|cccc}
      $x$ & $0$ & $1$ & $2$ & $3$ \\
      \hline
      $P(X=x)$ & $0.2$ & $0.3$ & $0.4$ & $0.1$ \\
    \end{tabular}
  \end{center}
  Tìm hàm phân phối xác suất của $X$.

  \begin{align*}
    F_X(x) = \begin{cases}
      0, & x < 0 \\
      0.2, & 0 \le x < 1 \\
      0.2 + 0.3 = 0.5, & 1 \le x < 2 \\
      0.2 + 0.3 + 0.4 = 0.9, & 2 \le x < 3 \\
      1, & x \ge 3
    \end{cases}
  \end{align*}
\end{vidu}

\subsubsection{Hàm mật độ}

\begin{dinhnghia}
  Giả sử  $X$ là một biến ngẫu nhiên liên tục có hàm phân phối xác suất $F_X(x), x \in R$. Nếu tồn tại hàm $f_X(x)$ sao cho:
  \[
    F_X(x) = \int_{-\infty}^{x} f_X(t)\,dt \forall x \in \mathbb{R} 
  \]
  thì $f_X(x)$ được gọi là hàm mật độ xác suất của biến ngẫu nhiên liên tục $X$.

  Hàm mật độ xác suất của biến ngẫu nhiên liên tục $X$ là đạo hàm bậc nhất của hàm phân phối xác suất của biến ngẫu nhiên đó,
  \[
    f_X(x) = F_X'(x), \forall x \in \mathbb{R}
  \]
\end{dinhnghia}

\begin{tinhchat}[title={Tính chất của hàm mật độ}]
  \begin{itemize}
    \item $f_X(x) \ge 0$ với mọi $x \in \mathbb{R}$.
    \item $P(a < X < b) = P(a \le X < b) = P(a \le X \le b) = P(a < X \le b) = F_X(b) - F_X(a) = \int_a^b f_X(x)\,dx$
    \item $\int_{-\infty}^{+\infty} f_X(x)\,dx = 1$
  \end{itemize}
\end{tinhchat}

\begin{vidu}[title={Hàm mật độ của biến ngẫu nhiên}]
Gọi biến ngẫu nhiên liên tục $X$ là cường độ dòng điện đo được trong một sợi dây đồng mỏng tính bằng miliampe (mA). Giả sử $X$ nhận giá trị trong đoạn [0; 20 mA] và hàm mật độ xác suất của $X$ là
  \[
    f_X(x) = \begin{cases}
      0, & x \notin [0; 20] \\
      0.05, & x \in [0; 20]
    \end{cases}
  \]  
  Xác suất để phép đo đo được cường độ dòng điện nhỏ hơn 10 miliampe là bao nhiêu?

  \[
    P(X < 10) = \int_{\infty}^{10} f_X(x)\,dx = \int_{0}^{10} 0.05\,dx = 0.05 \times 10 = 0.5
  \]
\end{vidu}


\subsubsection{Phân phối của $Y = g(X)$}

\begin{phuongphap}
  Cho $X$ là một biến ngẫu nhiên và $Y = g(X)$, trong đó, $g(x)$ là một hàm số thực cho trước.
  
  \textbf{Bảng phân phối xác suất:} Nếu $X$ là biến ngẫu nhiên rời rạc, hàm xác suất hay bảng phân phối xác suất của $Y = g(X)$ được xác định trực tiếp từ biến ngẫu nhiên $X$ ban đầu.
  \textbf{Hàm mật độ xác suất:} Nếu $X$ là biến ngẫu nhiên liên tục và $g(x)$ là hàm số liên tục, để tìm hàm mật độ xác suất $f_Y(y)$ từ $Y = g(X)$ và hàm mật độ xác suất của $X$ ta sẽ thực hiện theo quy trình hai bước:
  \[
    F_Y(y) = P(Y < y) = P(g(X) < y)
  \]
  \[
    f_Y(y) = \frac{dF_Y(y)}{dy}
  \]
\end{phuongphap}

\begin{vidu}
Cho $Z$ là biến ngẫu nhiên có bảng phân phối xác suất:
  \begin{center}
    \begin{tabular}{c|cccc}
      $z$ & 1 & 2 & 3 & 4 \\
      \hline
      $P(Z=z_i)$ & 0.2 & 0.3 & $p$ & 0.15 \\
    \end{tabular}
  \end{center}
  \begin{itemize}
    \item Tìm $p$.
    \item Tính $P(Z \ge 2)$.
    \item Lập bảng phân phối xác suất của biến ngẫu nhiên $Y = 5Z + 2$.
    \item Tìm hàm phân phối xác suất của $Y$.
  \end{itemize}
$Z$ rời rạc. Từ $\sum p_i = 1$ suy ra $p = 1 - 0.2 - 0.3 - 0.15 = 0.35$.

  \[
    P(Z \ge 2) = P(Z = 2) + P(Z = 3) + P(Z = 4) = 0.3 + 0.35 + 0.15 = 0.8
  \]
  \begin{center}
    \begin{tabular}{c|cccc}
      $y$ & 7 & 12 & 17 & 22 \\
      \hline
      $P(Y=y)$ & 0.2 & 0.3 & 0.35 & 0.15 \\
    \end{tabular}
  \end{center}
  \begin{align*}
    F_Y(y) = P(Y < y) = P(5Z + 2 < y) = P(Z < \frac{y-2}{5})
  \end{align*}
  \[
    f_Y(y) = \frac{dF_Y(y)}{dy} = \frac{1}{5} f_Z(\frac{y-2}{5})
  \]
  \[
    f_Y(y) = \frac{1}{5} \begin{cases}
      0, & y < 2 \\
      0.2, & 2 \le y < 7 \\
      0.3, & 7 \le y < 12 \\
      0.35, & 12 \le y < 17 \\
      0.15, & 17 \le y < 22 \\
      0, & y \ge 22 \\
    \end{cases}
  \]
\end{vidu}

\subsection{Kỳ vọng, phương sai, mốt, trung vị}

\begin{dongco}
Bảng phân phối liệt kê \emph{mọi} khả năng - đôi khi dài và khó nuốt. Hai số tóm tắt hữu ích nhất:

\textbf{Kỳ vọng $E(X)$:} ``trung bình có trọng số''. Ví dụ: trò chơi nhận $+10$ triệu với xác suất $0{,}01$, mất $1$ triệu với $0{,}99$ $\Rightarrow$ $E = 0{,}01 \times 10 - 0{,}99 \times 1 = -0{,}89$ triệu (về lâu dài thua). Nhà cái thiết kế game sao cho $E$ có lợi cho họ.

Kỳ vọng của một biến ngẫu nhiên $X$ phản ánh giá trị trung tâm của phân phối xác suất của biến ngẫu nhiên. Tuy nhiên, kỳ vọng không đưa ra một mô tả đầy đủ về hình dạng của phân phối.

\textbf{Phương sai $\mathrm{V}(X)$:} ``độ lan'' quanh $E(X)$. Hai phân phối cùng kỳ vọng nhưng một tập trung, một phân tán rộng - rủi ro khác nhau dù ``trung bình'' giống nhau.

Trong kỹ thuật phương sai đặc trưng cho mức độ phân tán của các chi tiết gia công hay sai số của thiết bị. Trong quản lý và kinh doanh thì phương sai đặc trưng cho mức độ rủi ro của các quyết định và thay cho “phương sai” người ta thường dùng thuật ngữ “độ biến động”. Như vậy, phương sai càng nhỏ thì chi tiết máy có độ chính xác càng cao, phương án đầu tư có độ rủi ro càng thấp.

\textbf{Độ lệch chuẩn $\sqrt{\mathrm{V}(X)}$:} Khi cần đánh giá mức độ phân tán của biến ngẫu nhiên theo đơn vị đo của nó người ta dùng độ lệch chuẩn vì độ lệch chuẩn có cùng đơn vị đo với đơn vị đo của biến ngẫu nhiên.
\end{dongco}

\begin{dinhnghia}[title={Kỳ vọng $E(X)$}]
  \begin{itemize}
    \item Rời rạc (hữu hạn giá trị): $E(X) = \sum_i x_i p_i$.
    \item Liên tục: $E(X) = \displaystyle\int_{-\infty}^{+\infty} x f_X(x)\,dx$ (nếu hội tụ).
    \item $E[g(X)]$: thay $x$ bằng $g(x)$ trong công thức trên.
  \end{itemize}
\end{dinhnghia}

\begin{vidu}
Trong ví dụ về bít bị lỗi trong 4 bít được truyền đi trước đó: $E(X) = 0 \cdot 0{,}6561 + 1 \cdot 0{,}2916 + 2 \cdot 0{,}0486 + 3 \cdot 0{,}0036 + 4 \cdot 0{,}0001 = 0{,}4$.

Trong ví dụ về cường độ dòng điện đo được trên một sợi dây đồng mỏng trước đó: $E(x) = \int_{0}^{20} x \cdot 0.05\,dx = 10$.
\end{vidu}

\begin{dinhnghia}[title={Kỳ vọng $E[g(X)]$}]
  \begin{itemize}
    \item Rời rạc (hữu hạn giá trị): $E[g(X)] = \sum_i g(x_i) p_i$.
    \item Liên tục: $E[g(X)] = \int_{-\infty}^{+\infty} g(x) f_X(x)\,dx$ (nếu hội tụ).
  \end{itemize}
\end{dinhnghia}

\begin{vidu}
Trong ví dụ về bít bị lỗi trong 4 bít được truyền đi trước đó, với $g(x) = x^2$: $E[g(X)] = 0^2 \cdot 0{,}6561 + 1^2 \cdot 0{,}2916 + 2^2 \cdot 0{,}0486 + 3^2 \cdot 0{,}0036 + 4^2 \cdot 0{,}0001 = 0{,}52$.

Trong ví dụ về cường độ dòng điện đo được trên một sợi dây đồng mỏng trước đó, với $g(x) = x^2$: $E[g(X)] = \int_{0}^{20} x^2 \cdot 0.05\,dx = 133.33$.
\end{vidu}

\begin{tinhchat}[title={Tính chất kỳ vọng}]
  $E(aX + b) = aE(X) + b$
  
  Nếu $X,Y$ độc lập thì $E(XY) = E(X)E(Y)$.

  $E[g(X) \pm h(X)] = E[g(X)] \pm E[h(X)]$.

  Hệ quả: $E(c) = c, E(cX) = cE(X)$.
\end{tinhchat}

\begin{dinhnghia}[title={Phương sai $\mathrm{V}(X)$ và độ lệch chuẩn $\sqrt{\mathrm{V}(X)}$}]
  Phương sai: $\sigma_X^2 = \mathrm{V}(X) = E\bigl[(X - E(X))^2\bigr] = E(X^2) - [E(X)]^2$
  \begin{itemize}
    \item Rời rạc (hữu hạn giá trị): $\sigma_X^2 = \mathrm{V}(X) = \sum_i (x_i - E(X))^2 p_i = \sum_i x_i^2 p_i - (\sum_i x_i p_i)^2$.
    \item Liên tục: $\sigma_X^2 = \mathrm{V}(X) = \int_{-\infty}^{+\infty} (x - E(X))^2 f_X(x)\,dx = \int_{-\infty}^{+\infty} x^2 f_X(x)\,dx - (E(X))^2 = \int_{-\infty}^{+\infty} x^2 f_X(x)\,dx - (\int_{-\infty}^{+\infty} x f_X(x)\,dx)^2$  (nếu hội tụ).
  \end{itemize}
  
  Độ lệch chuẩn: $\sigma_X = \sqrt{\sigma_X^2} = \sqrt{V(X)}$.
\end{dinhnghia}

\begin{vidu}
Trong ví dụ về bít bị lỗi trong 4 bít được truyền đi trước đó: $\sigma_X^2 = \mathrm{V}(X) = \sum_i (x_i - 0.4)^2 p_i = 0.36$, $\sigma_X = \sqrt{0.36} = 0.6$.

Trong ví dụ về cường độ dòng điện đo được trên một sợi dây đồng mỏng trước đó: $\sigma_X^2 = \mathrm{V}(X) = \int_{-\infty}^{+\infty} (x - 10)^2 \cdot 0.05\,dx = 33.33$.
  $\sigma_X = \sqrt{33.33} = 5.77$.
\end{vidu}

\begin{dinhnghia}[title={Phương sai $\mathrm{V}[g(X)]$ và độ lệch chuẩn $\sqrt{\mathrm{V}[g(X)]}$}]
  Phương sai: $\sigma_{g(X)}^2 = \mathrm{V}[g(X)] = E\bigl[(g(X) - E[g(X)])^2\bigr] = E(g^2(X)) - [E(g(X))]^2$
  \begin{itemize}
    \item Rời rạc (hữu hạn giá trị): $\sigma_{g(X)}^2 = \mathrm{V}[g(X)] = \sum_i (g(x_i) - E[g(X)])^2 p_i = \sum_i g^2(x_i) p_i - (\sum_i g(x_i) p_i)^2$.
    \item Liên tục: $\sigma_{g(X)}^2 = \mathrm{V}[g(X)] = \int_{-\infty}^{+\infty} (g(x) - E[g(X)])^2 f_X(x)\,dx = \int_{-\infty}^{+\infty} g^2(x) f_X(x)\,dx - (E[g(X)])^2 = \int_{-\infty}^{+\infty} g^2(x) f_X(x)\,dx - (\int_{-\infty}^{+\infty} g(x) f_X(x)\,dx)^2$  (nếu hội tụ).
  \end{itemize}
  
  Độ lệch chuẩn: $\sigma_{g(X)} = \sqrt{\sigma_{g(X)}^2} = \sqrt{\mathrm{V}[g(X)]}$.
\end{dinhnghia}

\begin{vidu}
Trong ví dụ về bít bị lỗi trong 4 bít được truyền đi trước đó, với $g(x) = x^2$: $\sigma_{g(X)}^2 = \mathrm{V}[g(X)] = \sum_i (x_i^2 - 0.52)^2 p_i = 0.08$, $\sigma_{g(X)} = \sqrt{0.08} = 0.28$.

Trong ví dụ về cường độ dòng điện đo được trên một sợi dây đồng mỏng trước đó, với $g(x) = x^2$: $\sigma_{g(X)}^2 = \mathrm{V}[g(X)] = \int_{-\infty}^{+\infty} (x^2 - 133.33)^2 \cdot 0.05\,dx = 1111.11$, $\sigma_{g(X)} = \sqrt{1111.11} = 33.33$.
\end{vidu}

\begin{tinhchat}[title={Tính chất phương sai}]
  $\mathrm{V}(aX + b) = a^2 \mathrm{V}(X)$
  
  $\mathrm{V}(X + Y) = \mathrm{V}(X) + \mathrm{V}(Y) + 2\mathrm{Cov}(X,Y)$ (Cov xét ở Phần sau).

  Hệ quả: $\mathrm{V}(c) = 0, \mathrm{V}(cX) = c^2 \mathrm{V}(X)$.
\end{tinhchat}

\begin{dinhnghia}[title={Mốt và trung vị}]
  \textbf{Mốt:} Mốt là giá trị của biến ngẫu nhiên $X$, ký hiệu là $mod(X)$, mà tại đó hàm mật độ xác suất $f_X(x)$ đạt cực đại (địa phương).

  Nếu $X$ là biến ngẫu nhiên rời rạc, thì $mod(X)$ là giá trị tại đó xác suất để $X = mod(X)$ là lớn nhất.

  \textbf{Trung vị:} Trung vị của biến ngẫu nhiên $X$ là giá trị của biến ngẫu nhiên $X$, ký hiệu là $med(X)$, mà tại đó $F_X(med(X)) = 1/2$.

  Phân vị mức $\alpha$ là giá trị của biến ngẫu nhiên $X$, ký hiệu là $x_\alpha$, sao cho $F_X(x_\alpha) = \alpha$.
\end{dinhnghia}

\begin{vidu}
Cho phân phối xác suất của $X$ là
  \begin{center}
    \begin{tabular}{c|ccc}
      $x$ & -1 & 0 & 1 \\
      \hline
      $P(X = x)$ & 1/6 & 2/3 & 1/6 \\
    \end{tabular}
  \end{center}
  Ta có, $mod(X) = 0$
\end{vidu}

\begin{vidu}
  Cho hàm mật độ xác suất của biến ngẫu nhiên liên tục $X$:
  \[
    f_X(x) = \begin{cases} 0, & x \le 0 \\ \frac{3}{4} x (2 - x), & 0 < x \le 2 \\ 0, & x > 2 \end{cases}
  \]

  Tìm $mod(X)$ và $med(X)$:

  $med(X)$ là nghiệm của phương trình $F_X(x) = 1/2 \Leftrightarrow \int_{-\infty}^{x} f_X(t)\,dt = 1/2 \Leftrightarrow \int_{-\infty}^{x} \frac{3}{4} t (2 - t)\,dt = 1/2 \Leftrightarrow x^3 - 3x^2 + 2 = 0$ với $0 < x \leq 2$.
  
  Suy ra $med(X) = 1$.

  Ta có:
  
  \[
    f'_X(x) = \begin{cases} 0, & x \le 0 \\ \frac{3}{2} (1 - x), & 0 < x < 2 \\ 0, & x > 2 \end{cases}
  \]

  Đổi dấu từ dương sang âm khi đi qua $x = 1$, do đó đạt cực đại tại điểm này, nên $mod(X) = 1$.

  Vậy $mod(X) = 1$ và $med(X) = 1$.
\end{vidu}

\subsection{Một số phân phối xác suất thông dụng}

\begin{dongco}
Không có một công thức ``vạn năng''. Mỗi mô hình mô tả một \emph{kiểu} ngẫu nhiên:
\begin{itemize}
  \item \textbf{Nhị thức:} đếm số lần \emph{thành công} trong $n$ lần thử độc lập (bắn cung, bít lỗi).
  \item \textbf{Poisson:} đếm sự kiện \emph{hiếm} trong khoảng thời gian (cuộc gọi đến tổng đài, lỗi trên km dây).
  \item \textbf{Chuẩn:} sai số đo, chiều cao tập thể, và (khi $n$ lớn) xấp xỉ nhị thức.
\end{itemize}
Chọn đúng mô hình $\Leftrightarrow$ chọn đúng câu hỏi thực tế.
\end{dongco}

\begin{congthuc}[title={Bảng tóm tắt}]
\small
\begin{center}
\renewcommand{\arraystretch}{1.2}
\begin{tabular}{llll}
\toprule
Tên & Ký hiệu & $P(X=k)$ hoặc $f(x)$ & $E(X)$, $\mathrm{V}(X)$ \\
\midrule
Đều rời rạc & $U\{x_1,\ldots,x_n\}$ & $1/n$ & $\bar{x}$, $\mathrm{V}(X)$ theo bảng \\
Bernoulli & $B(p)$ & $P(1)=p, P(0) = 1 - p$ & $p$, $p(1-p)$ \\
Nhị thức & $B(n,p)$ & $\binom{n}{k}p^k(1-p)^{n-k}$ & $np$, $np(1-p)$ \\
Poisson & $P(\lambda)$ & $e^{-\lambda}\lambda^k/k!$ & $\lambda$, $\lambda$ \\
Đều liên tục & $U([a,b])$ & $1/(b-a), a \le x \le b$ & $(a+b)/2$, $(b-a)^2/12$ \\
Mũ & $Exp(\lambda)$ & $\lambda e^{-\lambda x}$, $x\ge 0$ & $1/\lambda$, $1/\lambda^2$ \\
Chuẩn (Gauss) & $N(\mu,\sigma^2)$ & $\dfrac{1}{\sigma\sqrt{2\pi}}e^{-\dfrac{(x-\mu)^2}{2\sigma^2}}$ & $\mu$, $\sigma^2$ \\
Chuẩn tắc & $Z \sim N(0,1)$ & $\dfrac{1}{\sqrt{2\pi}}e^{-\dfrac{z^2}{2}}$ & $0$, $1$ \\
Khi-bình phương & $\chi^2(n)$ & $\dfrac{1}{2^{n/2}\Gamma(n/2)}x^{n/2-1}e^{-x/2}$ & $n$, $2n$ \\
Student & $t(n)$ & $\dfrac{\Gamma((n+1)/2)}{\sqrt{n\pi}\Gamma(n/2)}(1+\dfrac{t^2}{n})^{-(n+1)/2}$ & $0$, $\dfrac{n}{n-2}$ \\
Fisher & $F(m,n)$ & $\dfrac{\Gamma((m+n)/2)}{\Gamma(m/2)\Gamma(n/2)}\dfrac{m^{m/2}n^{n/2}x^{m/2-1}}{(mx+n)^{(m+n)/2}}$ & $\dfrac{n}{n-2}$ \\
\bottomrule
\end{tabular}
\end{center}
\normalsize
Với $\Gamma(n) = (n-1)!$ nếu $n$ là số nguyên dương, $\Gamma(n) = \int_0^\infty x^{n-1}e^{-x}\,dx$ nếu $n$ không phải là số nguyên dương.
\end{congthuc}

\subsubsection{Phân phối đều rời rạc}

\begin{dinhnghia}
  BNN rời rạc $X$ có \textbf{phân phối đều} tham số $n$ nếu nhận các giá trị $x_1, x_2, \ldots, x_n$ với xác suất bằng nhau:
  \[
    P(X = x_i) = \frac{1}{n}, \qquad i = 1, 2, \ldots, n.
  \]
  Ký hiệu: $X \sim U\{x_1, \ldots, x_n\}$ hoặc ``đều trên $n$ giá trị''.
\end{dinhnghia}

\begin{dinhly}
  Nếu $X$ đều trên các số nguyên liên tiếp $a, a+1, \ldots, b$ ($a < b$) thì
  \[
    E(X) = \frac{a+b}{2}, \qquad V(X) = \frac{(b-a+1)^2 - 1}{12}.
  \]
\end{dinhly}

\begin{vidu}[title={Chữ số đầu số sê-ri}]
$X$ = chữ số đầu của số sê-ri, $X \in \{0,1,\ldots,9\}$, $n = 10$:
$p_X(x) = 0{,}1$ với $x = 0,\ldots,9$; $E(X) = 4{,}5$; $V(X) = 8{,}25$.
\end{vidu}

\begin{vidu}[title={Số đường truyền đang dùng}]
Hệ thống 48 đường; $X$ = số đường đang sử dụng tại một thời điểm, giả sử $X$ đều trên $\{0,1,\ldots,48\}$:
\[
  E(X) = \frac{0+48}{2} = 24, \qquad V(X) = \frac{49^2 - 1}{12} = 200, \qquad \sigma_X \approx 14{,}14.
\]
Trung bình 24 đường nhưng độ phân tán lớn --- nhiều lúc xa 24.
\end{vidu}

\subsubsection{Phân phối Bernoulli}

\begin{dinhnghia}
  $X \sim B(p)$ (Bernoulli tham số $p \in (0,1)$) nếu $X \in \{0, 1\}$ và
  \[
    p_X(x) = \begin{cases}
      1-p, & x = 0, \\
      p,   & x = 1.
    \end{cases}
  \]
  Một phép thử Bernoulli: $X = 1$ nếu sự kiện $A$ xảy ra ($P(A) = p$), $X = 0$ nếu không.
\end{dinhnghia}

\begin{dinhly}
  $E(X) = p$, \quad $V(X) = p(1-p)$.
\end{dinhly}

\subsubsection{Phân phối nhị thức}

\begin{dinhnghia}
  $X \sim B(n; p)$ (nhị thức) nếu $X \in \{0, 1, \ldots, n\}$ và
  \[
    p_X(x) = \binom{n}{x} p^x (1-p)^{n-x}, \quad x = 0, 1, \ldots, n.
  \]
  Xuất phát từ khai triển $(p + q)^n$ với $q = 1-p$: $\sum_{k=0}^{n} \binom{n}{k} p^k q^{n-k} = 1$.

  \textbf{Mô hình:} $n$ phép thử Bernoulli \emph{độc lập}, cùng $p$; $X$ = số lần $A$ xảy ra.
\end{dinhnghia}

\begin{dinhly}
  \begin{itemize}
    \item $E(X) = np$, \quad $V(X) = np(1-p)$.
    \item Mốt: $(n+1)p - 1 \le \mathrm{mod}(X) \le (n+1)p$.
  \end{itemize}
\end{dinhly}

\begin{vidu}[title={Kiểm tra 20 sản phẩm}]
Tỷ lệ loại A = 50\%; lấy $n = 20$ sản phẩm, $X$ = số A $\Rightarrow$ $X \sim B(20; 0{,}5)$.
\[
  P(X = 5) = \binom{20}{5}(0{,}5)^5(0{,}5)^{15} \approx 0{,}0148.
\]
Với $p = 0{,}5$: $(n+1)p = 10{,}5$ $\Rightarrow$ $\mathrm{mod}(X) = 10$.
\end{vidu}

\begin{vidu}[title={Bít lỗi}]
$X \sim B(4; 0{,}1)$: $E(X) = 0{,}4$, $V(X) = 0{,}36$ (khớp tính trực tiếp từ bảng phân phối).
\end{vidu}

\subsubsection{Phân phối Poisson}

\begin{dinhnghia}
  $X \sim P(\lambda)$, $\lambda > 0$, nếu
  \[
    p_X(x) = \frac{e^{-\lambda}\lambda^x}{x!}, \quad x = 0, 1, 2, \ldots
  \]
  (dùng $e^\lambda = \sum_{x=0}^{\infty} \lambda^x/x!$ để có $\sum_x p_X(x) = 1$).
\end{dinhnghia}

\begin{ynghia}
  Trong thực tế với một số giả thiết thích hợp thì các biến ngẫu nhiên là các quá trình đếm sau: số cuộc gọi đến một tổng đài; số khách hàng đến một điểm phục vụ; số xe cộ qua một ngã tư; số tai nạn (xe cộ); số các sự cố xảy ra ở một địa điểm . . . trong một khoảng thời gian xác định nào đó sẽ có phân phối Poisson với tham số $\lambda$, trong đó $\lambda$ là giá trị trung bình diễn ra trong khoảng thời gian này.
\end{ynghia}

\begin{dinhly}
  $E(X) = \lambda$, \quad $V(X) = \lambda$.
\end{dinhly}

\begin{vidu}[title={Cuộc gọi tổng đài}]
Trung bình 2 cuộc/phút $\Rightarrow$ trong $2{,}5$ phút: $X \sim P(5)$ với $\lambda = 2 \times 2{,}5 = 5$,
\[
  P(X = 6) = e^{-5}\frac{5^6}{6!} \approx 0{,}1462.
\]
Trong 30 giây: $Y \sim P(1)$ với $\lambda = 2 \times 0{,}5 = 1$, $P(Y = 0) = e^{-1} \approx 0{,}3679$.

Trong 10 giây: $Z \sim P(1/3)$ với $\lambda = 2 \times 1/3 = 2/3$, $P(Z \ge 1) = 1 - P(Z = 0) = 1 - e^{-1/3} \approx 0{,}2835$.
\end{vidu}

\begin{dinhly}[title={Xấp xỉ nhị thức bằng Poisson}]
  Cho $X$ là biến ngẫu nhiên có phân phối nhị thức với tham số $n$ và $p$. Nếu $n \rightarrow \infty$, $p \rightarrow 0$ và $np \rightarrow \lambda$ ($\lambda$ là một hằng số) thì
  \[
    B(n; p) \rightarrow P(\lambda)
  \]

  Trong thực tế, nếu $n$ đủ lớn và $p$ đủ nhỏ thỏa mãn $np < 7$ thì ta có thể xấp xỉ phân phối nhị thức $B(n; p)$ bằng phân phối Poisson $P(\lambda)$ với $\lambda = np$ và
  \[
    P_n(k) = C_k^{np} (1 - p)^{n-k} \approx \dfrac{(\lambda)^k}{k!} e^{-\lambda}
  \]
\end{dinhly}

\begin{vidu}[title={Bảo hiểm 5000 người}]
Giả sử một công ty bảo hiểm nhân thọ bảo hiểm cho cuộc sống của 5000 người đàn ông ở độ tuổi 42. Nghiên cứu của các chuyên gia tính toán cho thấy xác suất để một người đàn ông 42 tuổi sẽ chết trong một năm (xác định) là 0,001. Hãy tìm xác suất mà công ty sẽ phải trả bảo hiểm cho 4 người trong một năm (xác định).

$n = 5000$, $p = 0{,}001$, $\lambda = 5 < 7$. $P(X = 4)$:

Công thức Bernoulli: $\approx 0{,}1755$;

Xấp xỉ Poisson $\approx 0{,}175$ (gần nhau).
\end{vidu}

\subsubsection{Phân phối đều liên tục}

\begin{dinhnghia}
  $X \sim U([a,b])$, $a < b$, nếu
  \[
    f_X(x) = \begin{cases}
      \dfrac{1}{b-a}, & x \in [a,b], \\
      0, & \text{khác}.
    \end{cases}
  \]
\end{dinhnghia}

\begin{tinhchat}
  Với $\alpha, \beta \in [a,b]$: $P(\alpha < X < \beta) = \dfrac{\beta - \alpha}{b - a}$.
\end{tinhchat}

\begin{dinhly}
Hàm phân phối xác suất của $X$ là:
  \[
    F_X(x) = \begin{cases}
      0, & x \le a, \\
      \dfrac{x-a}{b-a}, & a < x \le b, \\
      1, & x > b.
    \end{cases}
  \]
  $E(X) = \dfrac{a+b}{2}$, \quad $V(X) = \dfrac{(b-a)^2}{12}$.
\end{dinhly}

\begin{vidu}[title={Chờ xe buýt}]
Xe mỗi 15 phút từ 7h00; khách đến ngẫu nhiên $X \sim U[0;30]$ (phút, quy 7h = 0).
\begin{itemize}
  \item Chờ $< 5$ phút: hai đoạn $(10,15]$ và $(25,30]$ $\Rightarrow$ $P = \dfrac{15-10}{30-0} + \dfrac{30-25}{30-0} = \dfrac{5+5}{30} = \dfrac{1}{3}$.
  \item Chờ $\ge 12$ phút: $(0,3] \cup (15,18]$ $\Rightarrow$ $P = \dfrac{3-0}{30-0} + \dfrac{18-15}{30-0} = \dfrac{3+3}{30} = \dfrac{1}{5}$.
\end{itemize}
\end{vidu}

\subsubsection{Phân phối mũ}

\begin{dinhnghia}
  $X \sim \mathrm{Exp}(\lambda)$, $\lambda > 0$, nếu
  \[
    f_X(x) = \begin{cases}
      \lambda e^{-\lambda x}, & x \ge 0, \\
      0, & x < 0.
    \end{cases}
  \]
  Mô hình \textbf{thời gian chờ} đến sự kiện đầu tiên (hỏng hóc, cuộc gọi tiếp theo\ldots).
\end{dinhnghia}

\begin{dinhly}
  \[
    F_X(x) = \begin{cases}
      1 - e^{-\lambda x}, & x \ge 0, \\
      0, & x < 0.
    \end{cases}
  \]
  $E(X) = 1/\lambda$, \quad $V(X) = 1/\lambda^2$.

  \textbf{Hệ quả:} $P(X > x) = e^{-\lambda x}$, $x \ge 0$.
\end{dinhly}

\begin{vidu}[title={Thời gian $T$}]
Biến ngẫu nhiên $T$ có phân phối mũ với tham số $\lambda > 0$ có hàm phân phối xác suất là
  \[
    F_T(t) = \begin{cases}
      1 - e^{-t/3}, & t \ge 0, \\
      0, & t < 0.
    \end{cases}
  \]
  \begin{itemize}
    \item Tìm hàm mật độ xác suất của $T$.
    \item Tính $P(2 \le T \le 4)$.
    \item Tính kỳ vọng, phương sai, độ lệch chuẩn của $T$.
  \end{itemize}
$f_T(t) = F_T'(t) = \frac{1}{3}e^{-t/3}$; 

$P(2 \le T \le 4) = e^{-2/3} - e^{-4/3} \approx 0{,}250$;

$E(T) = 3$, $V(T) = 9$, $\sigma_T = 3$.
\end{vidu}

\begin{vidu}[title={Tuổi thọ mạch điện tử}]
Giả sử tuổi thọ (năm) của một mạch điện tử trong máy tính là một biến ngẫu nhiên có phân phối mũ với kỳ vọng là 6,25 năm. Thời gian bảo hành của mạch điện tử này là 5 năm. Hỏi có bao nhiêu phần trăm mạch điện tử bán ra phải thay thế trong thời gian bảo hành.

$E(X) = 6{,}25$ năm $\Rightarrow$ $\lambda = 1/6{,}25 = 0{,}16$.

Bảo hành 5 năm: $P(X \le 5) = 1 - e^{-0{,}16 \times 5} = 1 - e^{-0{,}8} \approx 55{,}06\%$ phải thay trong bảo hành.
\end{vidu}

\begin{tinhchat}[title={Tính không nhớ}]
  $X \sim \mathrm{Exp}(\lambda)$: với mọi $s, t \ge 0$,
  \[
    P(X > s + t \mid X > t) = P(X > s).
  \]
\end{tinhchat}

\begin{vidu}[title={Thời gian chờ đến sự kiện đầu tiên}]
Một động cơ có thời gian làm việc (giờ) liên tục cho đến khi có hỏng hóc đầu tiên được giả thiết là biến ngẫu nhiên có phân phối mũ với tham số λ = 0, 001.
\begin{itemize}
  \item Tính xác suất để động cơ này làm việc liên tục không hỏng hóc được ít nhất 900 giờ.
  \item Được biết động cơ này đã làm việc liên tục không hỏng hóc được ít nhất 300 giờ, xác suất để động cơ còn làm việc liên tục không hỏng hóc được ít nhất 1200 giờ là bao nhiêu?
\end{itemize}
Gọi $X$ là thời gian làm việc của động cơ, $X \sim \mathrm{Exp}(\lambda)$ với $\lambda = 0{,}001$.
\[
  P(X > 900) = e^{-\lambda \times 900} = e^{-0{,}9} \approx 0{,}4066.
\]
Áp dụng công thức xác suất điều kiện,
\[
  P(X > 1200 \mid X > 300) = \frac{P(X > 1200)}{P(X > 300)} = \frac{e^{-\lambda \times 1200}}{e^{-\lambda \times 300}} = e^{-\lambda (1200-300)} \approx 0{,}4066.
\]
\end{vidu}

\subsubsection{Phân phối chuẩn (Gauss)}

\begin{dinhnghia}
  $X \sim N(\mu; \sigma^2)$, $\sigma > 0$, nếu
  \[
    f_X(x) = \frac{1}{\sigma\sqrt{2\pi}}e^{-\frac{(x-\mu)^2}{2\sigma^2}}, \quad x \in \mathbb{R}.
  \]
  Đồ thị $f_X$ là \textbf{đường cong chuẩn} (hình chuông).
  
  \textbf{Ứng dụng:} chiều cao, trọng lượng, sai số đo, điểm thi (khi tập thể lớn)\ldots
\end{dinhnghia}

\begin{dinhly}
  $E(X) = \mu$, \quad $V(X) = \sigma^2$, \quad $\sigma_X = \sigma$.

  \textbf{Biến đổi tuyến tính:} $Y = \alpha X + \beta \Rightarrow Y \sim N(\alpha\mu + \beta;\, \alpha^2\sigma^2)$.
\end{dinhly}

\paragraph{Chuẩn tắc và hàm phân phối.}

\begin{dinhnghia}
  $Z = \dfrac{X - \mu}{\sigma}$ với $X \sim N(\mu;\sigma^2) \Rightarrow E[Z] = 0, V[Z] = 1 \Rightarrow Z \sim N(0;1)$ (chuẩn tắc).

  Hàm mật độ chuẩn tắc:
  \[
    \varphi(z) = \dfrac{1}{\sqrt{2\pi}} e^{-z^2/2}.
  \]

  Hàm phân phối chuẩn tắc:
  \[
    \Phi(z) = P(Z < z) = \dfrac{1}{\sqrt{2\pi}} \int_{-\infty}^{z} e^{-u^2/2}\,du.
  \]
\end{dinhnghia}

\begin{dinhly}
  $\Phi(-z) = 1 - \Phi(z)$.

  Với $X \sim N(\mu;\sigma^2)$:
  \[
    F_X(x) = \Phi\!\left(\frac{x-\mu}{\sigma}\right), \qquad
  \]
  \[
    P(\alpha < X < \beta) = F_X(\beta) - F_X(\alpha) = \Phi\!\left(\frac{\beta-\mu}{\sigma}\right) - \Phi\!\left(\frac{\alpha-\mu}{\sigma}\right).
  \]

  Hệ quả:
  \[
    P(X < \beta) = \Phi\!\left(\frac{\beta-\mu}{\sigma}\right).
  \]
  \[
    P(X > \alpha) = 1 - \Phi\!\left(\frac{\alpha-\mu}{\sigma}\right).
  \]
  \[
    P(|X - \mu| < \epsilon) = 2\Phi\!\left(\frac{\epsilon}{\sigma}\right) - 1.
  \]
\end{dinhly}

\begin{vidu}[title={Đo dòng điện}]
Giả sử phép đo cường độ dòng điện trên một sợi dây đồng là biến ngẫu nhiên có phân phối chuẩn với giá trị trung bình là 10 miliampe và phương sai là 4 (miliampe)$^2$. Tính xác suất để nhận được kết quả của phép đo này vượt quá 13 miliampe.

$X \sim N(10; 4)$ (mA), $\sigma = 2$:

$P(X > 13) = 1 - \Phi\!\left(\frac{13-10}{2}\right) = 1 - \Phi(1{,}5) = 1 - 0{,}93319 \approx 0{,}06681$.
\end{vidu}

\begin{vidu}[title={Lãi suất đầu tư}]
Lãi suất (\%) đầu tư vào một dự án trong năm 2022 được giả thiết là biến ngẫu nhiên có phân phối chuẩn. Theo đánh giá của ủy ban đầu tư thì với xác suất 0,15866 cho lãi suất lớn hơn 20\% và với xác suất 0,02275 cho lãi suất lớn hơn 25\%. Vậy khả năng đầu tư mà không bị lỗ là bao nhiêu?

$P(X > 20) = 1 - \Phi\!\left(\frac{20-\mu}{\sigma}\right) = 0{,}15866$

$P(X > 25) = 1 - \Phi\!\left(\frac{25-\mu}{\sigma}\right) = 0{,}02275$

\Phi\!\left(\frac{20-\mu}{\sigma}\right) = 1 - 0{,}15866 = 0{,}84134

$\Phi\!\left(\frac{25-\mu}{\sigma}\right) = 1 - 0{,}02275 = 0{,}97725$

$\Rightarrow$ từ bảng $\Phi$ suy ra $\frac{20-\mu}{\sigma} = 1$ và $\frac{25-\mu}{\sigma} = 2 \Rightarrow \mu = 15\%$ và $\sigma = 5\%$.

Không lỗ: $P(X \ge 0) = 1 - \Phi\!\left(\frac{0-\mu}{\sigma}\right) = 1 - \Phi(-3) = \Phi(3) \approx 0{,}99865$.
\end{vidu}

\begin{luuy}[title={Nhận xét}]
  $X \sim N(\mu;\sigma^2)$:
  \begin{align*}
    P(\mu - \sigma < X < \mu + \sigma) &\approx 0{,}68268, \\
    P(\mu - 2\sigma < X < \mu + 2\sigma) &\approx 0{,}9545, \\
    P(\mu - 3\sigma < X < \mu + 3\sigma) &\approx 0{,}9973.
  \end{align*}
\end{luuy}

\begin{dinhly}[title={Moivre--Laplace}]
  $X \sim B(n;p)$. Nếu $np \ge 5$ và $n(1-p) \ge 5$ thì $X$ xấp xỉ $N(np;\, np(1-p))$.
\end{dinhly}

\begin{dinhly}[title={Xấp xỉ phân phối nhị thức bằng phân phối chuẩn}]
  \[
    P(X = k) = P(k \le X \le k) \approx \Phi\!\left(\frac{k + 0{,}5 - np}{\sqrt{np(1-p)}}\right) - \Phi\!\left(\frac{k - 0{,}5 - np}{\sqrt{np(1-p)}}\right),
  \]
  \[
    P(k_1 \le X \le k_2) \approx \Phi\!\left(\frac{k_2 + 0{,}5 - np}{\sqrt{np(1-p)}}\right) - \Phi\!\left(\frac{k_1 - 0{,}5 - np}{\sqrt{np(1-p)}}\right).
  \]
\end{dinhly}

\begin{vidu}[title={$B(25; 0{,}5)$, $P(8 \le X \le 10)$}]
Sử dụng phân phối chuẩn xấp xỉ xác suất $P(8 \le X \le 10)$ cho biến ngẫu nhiên $X \sim B(25; 0{,}5)$. So sánh với công thức tính chính xác.

Chính xác:
\begin{align*}
  P(8 \le X \le 10) = C_{25}^{8} (0{,}5)^{8} (0{,}5)^{17} + C_{25}^{9} (0{,}5)^{9} (0{,}5)^{16} + C_{25}^{10} (0{,}5)^{10} (0{,}5)^{15} \\ = C_{25}^{8} (0{,}5)^{25} + C_{25}^{9} (0{,}5)^{25} + C_{25}^{10} (0{,}5)^{25} = 0{,}190535.
\end{align*}

Xấp xỉ chuẩn:
\begin{align*}
  P(8 \le X \le 10) \approx \Phi\!\left(\frac{10 + 0{,}5 - 12{,}5}{\sqrt{12{,}5 \times 0{,}5}}\right) - \Phi\!\left(\frac{7 - 0{,}5 - 12{,}5}{\sqrt{12{,}5 \times 0{,}5}}\right) \\= \Phi(-0{,}8) - \Phi(-2) = \Phi(2) - \Phi(0{,}8) = 0.97725 - 0.78814 = 0{,}18911.
\end{align*}
\end{vidu}

\begin{vidu}[title={1000 sản phẩm, 95\% chính phẩm}]
Kiểm tra chất lượng 1000 sản phẩm trong một lô hàng có tỷ lệ chính phẩm là 0,95. Tìm xác suất để trong 1000 sản phẩm đó có từ 940 đến 960 chính phẩm

Gọi $X$ là biến ngẫu nhiên chỉ số chính phẩm trong 1000 sản phẩm được kiểm tra, $X \sim B(n; p)$ với $n = 1000$ và $p = 0{,}95$.

Vì $np = 950$ và $n(1 - p) = 50$ khá lớn, nên ta có thể xấp xỉ phân phối nhị thức $B(n; p)$ bởi phân phối chuẩn $N (\mu; \sigma^2)$ với $\mu = np = 950$, $\sigma^2 = np(1 - p) = 47{,}5$ và 

\begin{align*}
P(940 \le X \le 960) \approx \Phi\!\left(\frac{960 + 0{,}5 - 950}{\sqrt{47{,}5}}\right) - \Phi\!\left(\frac{940 - 0{,}5 - 950}{\sqrt{47{,}5}}\right) \\= \Phi(1{,}52) - \Phi(-1{,}52) = 2\Phi(1{,}52) - 1 = 2 \times 0{,}93574 - 1 = 0{,}87148
\end{align*}
\end{vidu}

\subsubsection{Phân phối Khi-bình phương}

\begin{dinhnghia}
  $X \sim \chi^2(\nu)$ ($\nu$ bậc tự do) nếu
  \[
    f_X(x) = \begin{cases}
      \dfrac{1}{2^{\nu/2}\Gamma(\nu/2)}\, x^{\nu/2 - 1} e^{-x/2}, & x > 0, \\
      0, & x \le 0,
    \end{cases}
  \]
  với $\Gamma(x) = \int_0^{\infty} t^{x-1} e^{-t}\,dt$ (nếu $x$ nguyên dương: $\Gamma(n) = (n-1)!$).
\end{dinhnghia}

\begin{ynghia}
  Đồ thị của phân phối Khi-bình phương bị lệch sang phải. Tuy nhiên, khi số bậc tự do $ν$ tăng lên, phân phối này trở nên đối xứng hơn. Khi $ν \to \infty$, dạng giới hạn của phân phối Khi-bình phương là phân phối chuẩn.
\end{ynghia}

\begin{dinhly}
  $E(X) = \nu$, \quad $V(X) = 2\nu$.
\end{dinhly}

\begin{dinhnghia}[title={Giá trị tới hạn $\chi^2_{\alpha;\nu}$}]
  Ký hiệu $\chi^2_{\alpha;\nu}$ chỉ giá trị của biến ngẫu nhiên $X \sim \chi^2(\nu)$ sao cho xác suất $X$ vượt quá giá trị này là $\alpha$.
  Nghĩa là, $P(X > \chi^2_{\alpha;\nu}) = \alpha$.

  Giá trị $\chi^2_{\alpha;\nu}$ được tính sẵn trong Bảng giá trị tới hạn phân phối Khi-bình phương.

  Ví dụ: $P(X > \chi^2_{0{,}025;9}) = P(X > 19{,}023) = 0{,}025$.

  Ngược lại, ta có $\chi^2_{0{,}975;9} = 2{,}700$.
\end{dinhnghia}

\subsubsection{Phân phối Student}

\begin{dinhnghia}
  $T \sim t(\nu)$ nếu mật độ có dạng:
  \[
    f_T(t) = \frac{\Gamma\bigl(\frac{\nu+1}{2}\bigr)}{\sqrt{\nu\pi}\,\Gamma\bigl(\frac{\nu}{2}\bigr)} \left(1 + \frac{t^2}{\nu}\right)^{-(\nu+1)/2}, \quad t \in \mathbb{R}.
  \]
  Dạng ``đuôi dày'' hơn chuẩn; khi $\nu \ge 30$ có thể xấp xỉ $N(0,1)$.
\end{dinhnghia}

\begin{dinhly}[title={Cấu tạo từ chuẩn và $\chi^2$ (Định lý 19)}]
  $Z \sim N(0,1)$, $V \sim \chi^2(\nu)$ độc lập $\Rightarrow$
  \[
    T = \frac{Z}{\sqrt{V/\nu}} \sim t(\nu).
  \]
\end{dinhly}

\begin{dinhly}
  $E(T) = 0$ ($\nu > 1$); \quad $V(T) = \dfrac{\nu}{\nu-2}$ ($\nu > 2$).
\end{dinhly}

\begin{dinhnghia}[title={Giá trị tới hạn $t_{\alpha;\nu}$}]
  Ký hiệu $t_{\alpha;\nu}$ chỉ giá trị của biến ngẫu nhiên $T \sim t(\nu)$ sao cho xác suất $T$ vượt quá giá trị này là $\alpha$. Nghĩa là,
  $P(T > t_{\alpha;\nu}) = \alpha$.

  Giá trị $t_{\alpha;\nu}$ được tính sẵn trong Bảng giá trị tới hạn phân phối Student.

  Chẳng hạn, giá trị với 10 bậc tự do có xác suất ở bên phải bằng 0,05 là $t_{0{,}05;10} = 1{,}812$, tức là
  $P(T > t_{0{,}05;10}) = P(T > 1{,}812) = 0{,}05$.

  Vì tính đối xứng nên $t_{0{,}95;10} = -t_{0{,}05;10} = -1{,}812$.
\end{dinhnghia}

\subsubsection{Phân phối Fisher}

\begin{dinhnghia}
  $X \sim F(\nu_1; \nu_2)$ nếu mật độ (slide):
  \[
    f_X(x) = \frac{\Gamma\bigl[(\nu_1+\nu_2)/2\bigr]}{\Gamma(\nu_1/2)\Gamma(\nu_2/2)} \frac{(\nu_1/\nu_2)^{\nu_1/2}\, x^{\nu_1/2 - 1}}{\bigl(1 + \nu_1 x / \nu_2\bigr)^{(\nu_1+\nu_2)/2}}, \quad x > 0.
  \]
  Phụ thuộc \emph{thứ tự} $(\nu_1, \nu_2)$.
\end{dinhnghia}

\begin{dinhly}[title={Cấu tạo}]
  $U \sim \chi^2(\nu_1)$, $V \sim \chi^2(\nu_2)$ độc lập $\Rightarrow$
  \[
    F = \frac{U/\nu_1}{V/\nu_2} \sim F(\nu_1; \nu_2).
  \]
\end{dinhly}

\begin{dinhly}
  $E(X) = \dfrac{\nu_2}{\nu_2 - 2}$ ($\nu_2 > 2$).

  $V(X) = \dfrac{2\nu_2^2(\nu_1 + \nu_2 - 2)}{\nu_1(\nu_2-2)^2(\nu_2-4)}$ ($\nu_2 > 4$).

  $F_{1-\alpha;\nu_1;\nu_2} = \frac{1}{F_{\alpha;\nu_2;\nu_1}}$.
\end{dinhly}

\begin{dinhnghia}[title={Giá trị tới hạn $F_{\alpha;\nu_1;\nu_2}$}]
  Ký hiệu $F_{\alpha;\nu_1;\nu_2}$ là giá trị tới hạn phân phối Fisher ứng với bậc tự do ν1 và ν2 mà ở đó ta tìm được diện tích của miền phẳng được giới hạn bởi đường cong $f_X(x)$, trục hoành và bên phải của $F_{\alpha;\nu_1;\nu_2}$ bằng α
  Nghĩa là, $P(X > F_{\alpha;\nu_1;\nu_2}) = \alpha$.

  Giá trị $F_{\alpha;\nu_1;\nu_2}$ được tính sẵn trong Bảng giá trị tới hạn phân phối Fisher.

  Ví dụ: $P(X > F_{0{,}05;6;10}) = P(X > 4{,}06) = 0{,}05$.

  Ngược lại, ta có $F_{0{,}95;6;10} = \frac{1}{F_{0{,}05;10;6}} = \frac{1}{4{,}06} = 0{,}2463$.
\end{dinhnghia}

%% ================================================================
\section{Biến ngẫu nhiên nhiều chiều}

\begin{dongco}[title={Động cơ học -- Chiều cao và cân nặng}]
Một BNN $X$ (chiều cao) cho biết ``cao hay thấp''. Nhưng hai người cùng $170$ cm có thể nặng $55$ kg hoặc $85$ kg - \textbf{một số không đủ}.

Ta cần $(X,Y)$ \emph{đồng thời}: bảng $p_{ij} = P(X=x_i, Y=y_j)$ hoặc mật độ $f_{X,Y}(x,y)$. Từ đó:
\begin{itemize}
  \item \textbf{Biên:} chỉ quan tâm $X$ (bỏ qua $Y$) - cộng theo $y$.
  \item \textbf{Có điều kiện:} đã biết $X = 170$, phân phối của $Y$ thay đổi thế nào?
  \item \textbf{Độc lập:} biết $X$ \emph{không} giúp đoán $Y$ hơn - $f_{X,Y} = f_X f_Y$.
\end{itemize}
Đây là nền cho hồi quy, tương quan, và mọi mô hình ``nhiều biến'' sau này.
\end{dongco}

\subsection{Khái niệm BNN nhiều chiều}

\begin{dinhnghia}
  \textbf{BNN $n$ chiều} $(X_1, \ldots, X_n)$: các $X_i$ xác định trên cùng một phép thử.
  \textbf{Hai chiều} $(X,Y)$ là trường hợp trung tâm của giáo trình.
\end{dinhnghia}

\begin{vidu}[title={Sản phẩm đúc}]
$(X,Y,Z)$: chiều rộng, dài, cao. Có thể xét $(X,Y)$ nếu bỏ qua chiều cao.
\end{vidu}

\subsection{Phân phối đồng thời, biên, có điều kiện}

\begin{dinhnghia}[title={Phân phối đồng thời}]
  \textbf{Rời rạc:} $p_{ij} = P(X = x_i, Y = y_j)$, $\sum_{i,j} p_{ij} = 1$ (bảng hai chiều).

  \textbf{Liên tục:} $f_{X,Y}(x,y) \ge 0$, $\iint f_{X,Y} = 1$; $F_{X,Y}(x,y) = P(X \le x, Y \le y)$.
\end{dinhnghia}

\begin{dinhnghia}[title={Phân phối biên và có điều kiện}]
  \[
    p_{i\cdot} = \sum_j p_{ij} = P(X = x_i), \qquad
    f_X(x) = \int f_{X,Y}(x,y)\,dy.
  \]
  \[
    P(Y = y_j \mid X = x_i) = \frac{p_{ij}}{p_{i\cdot}}, \qquad
    f_{Y|X}(y|x) = \frac{f_{X,Y}(x,y)}{f_X(x)} \quad (f_X(x) > 0).
  \]
\end{dinhnghia}

\begin{dinhnghia}[title={Độc lập}]
  $X$, $Y$ \textbf{độc lập} $\Leftrightarrow$ $p_{ij} = p_{i\cdot}\, p_{\cdot j}$ (rời rạc) hoặc $f_{X,Y}(x,y) = f_X(x) f_Y(y)$ (liên tục).

  \textbf{Độc lập từng đôi} $\not\Rightarrow$ \textbf{độc lập tổng thể}.
\end{dinhnghia}

\begin{congthuc}[title={Kỳ vọng của hàm hai BNN}]
  $E[g(X,Y)] = \sum_{i,j} g(x_i,y_j) p_{ij}$ \; (rời rạc); tương tự tích phân kép nếu liên tục.
\end{congthuc}

\begin{vidu}[title={Tung xu 3 lần: $X$ = sấp, $Y$ = ngửa}]
$X + Y = 3$. Với mỗi cặp $(x,y)$ thỏa $x+y \le 3$, $x,y \ge 0$:
\[
  P(X=x, Y=y) = \frac{\binom{5}{x}\binom{4}{y}\binom{3}{3-x-y}}{\binom{12}{3}}.
\]
Bảng đồng thời (một phần):

\small
\begin{center}
\begin{tabular}{c|cccc}
$X \backslash Y$ & 0 & 1 & 2 & 3 \\
\hline
0 & $1/220$ & $3/55$ & $9/110$ & $1/55$ \\
1 & $3/44$ & $3/11$ & $3/22$ & $0$ \\
2 & $3/22$ & $2/11$ & $0$ & $0$ \\
3 & $1/22$ & $0$ & $0$ & $0$ \\
\end{tabular}
\end{center}
\normalsize

$P(X+Y \le 1) = P(0,0)+P(0,1)+P(1,0) = \dfrac{1}{220}+\dfrac{3}{55}+\dfrac{3}{44} = \dfrac{7}{55}$.
\end{vidu}

\subsection{Hiệp phương sai và hệ số tương quan}

\begin{dinhnghia}
  $\mathrm{Cov}(X,Y) = E\bigl[(X - EX)(Y - EY)\bigr] = E(XY) - EX \cdot EY$.

  \textbf{Hệ số tương quan:} $\rho_{XY} = \dfrac{\mathrm{Cov}(X,Y)}{\sigma_X \sigma_Y}$, $|\rho| \le 1$.
  \begin{itemize}
    \item $\mathrm{Cov} > 0$: $Y$ có xu hướng tăng khi $X$ tăng.
    \item $\mathrm{Cov} < 0$: xu hướng ngược.
    \item $\mathrm{V}(X) = \mathrm{Cov}(X,X)$.
  \end{itemize}
\end{dinhnghia}

\begin{luuy}
Độc lập $\Rightarrow$ $\mathrm{Cov}(X,Y) = 0$ (không tương quan tuyến tính). Ngược lại \emph{không} đúng nói chung (ví dụ $Y = X^2$).
\end{luuy}

\begin{tongket}[title={Tổng kết}]
Đồng thời $\to$ biên $\to$ có điều kiện; độc lập; $\mathrm{Cov}$, $\rho$.
\end{tongket}

%% ================================================================
\section{Luật số lớn}

\begin{dongco}[title={Động cơ học -- Thăm dò ý kiến và đo lường}]
\textbf{Thăm dò 1000 cử tri} trong tổng số hàng triệu: tỷ lệ ủng hộ $52\% \pm$ sai số. Tại sao ta \emph{tin} con số này hơn hỏi đúng 3 người bạn?

\textbf{Gieo xu 10\,000 lần:} tần suất sấp gần $0{,}5$; gieo 10 lần mà ra 9 sấp vẫn có thể xảy ra - không chứng minh xu gian.

\textbf{Đo điện trở 50 lần:} mỗi lần có sai số ngẫu nhiên; lấy trung bình $\bar{X}_{50}$ ổn định hơn một lần đo.

Ba hiện tượng cùng một ý: \textbf{trung bình mẫu} (hoặc tần suất) \emph{hội tụ} tới giá trị ``thật'' khi $n \to \infty$ - nhưng cần nói rõ \emph{theo xác suất} (LSL) và dùng bất đẳng thức Chebyshev để ước lượng ``sai bao nhiêu với $n$ quan sát''.
\end{dongco}

\subsection{Hội tụ theo xác suất}

\begin{ynghia}
Dãy $Z_n$ \textbf{hội tụ theo xác suất} tới $c$ nếu $P(|Z_n - c| < \varepsilon) \to 1$ với mọi $\varepsilon > 0$ (ký hiệu $Z_n \xrightarrow{P} c$).
\end{ynghia}

\subsection{Bất đẳng thức Markov và Chebyshev}

\begin{dinhly}[title={Bất đẳng thức Markov}]
Nếu $Y \ge 0$ và $E(Y)$ hữu hạn thì với $a > 0$: $P(Y \ge a) \le \dfrac{E(Y)}{a}$.
\end{dinhly}

\begin{dinhly}[title={Bất đẳng thức Chebyshev}]
$E(X) = \mu$, $\mathrm{r}(X) = \sigma^2 < \infty$. Với $\varepsilon > 0$:
\[
  P(|X - \mu| \ge \varepsilon) \le \frac{\sigma^2}{\varepsilon^2}
  \quad\Leftrightarrow\quad
  P(|X - \mu| < \varepsilon) \ge 1 - \frac{\sigma^2}{\varepsilon^2}.
\]
Đặt $\varepsilon = k\sigma$: $P(\mu - k\sigma < X < \mu + k\sigma) \ge 1 - 1/k^2$ (với $k=2$: $\ge 3/4$; $k=3$: $\ge 8/9$).
\end{dinhly}

\begin{vidu}[title={Chebyshev khi chưa biết phân phối}]
$EX = 8$, $\mathrm{V}(X) = 9$, $\sigma = 3$.
\begin{itemize}
  \item $P(-4 < X < 20) = P(8 - 4\cdot 3 < X < 8 + 4\cdot 3) \ge 1 - 1/16 = 15/16$.
  \item $P(|X - 8| \ge 6) = P(|X - 8| \ge 2\sigma) \le 1/4$.
\end{itemize}
\end{vidu}

\subsection{Luật số lớn}

\begin{dinhly}[title={Luật số lớn Chebyshev}]
$X_1, X_2, \ldots$ độc lập, $EX_i$ hữu hạn, $\mathrm{V}(X_i) \le C$. Với $\varepsilon > 0$:
\[
  \lim_{n \to \infty} P\!\left(\left|\frac{1}{n}\sum_{i=1}^n X_i - \frac{1}{n}\sum_{i=1}^n EX_i\right| < \varepsilon\right) = 1.
\]
Nếu cùng $EX_i = \mu$: $\displaystyle\frac{1}{n}\sum_{i=1}^n X_i \xrightarrow{P} \mu$.
\end{dinhly}

\begin{dinhly}[title={Luật số lớn Bernoulli}]
$n$ phép thử Bernoulli, $P(A) = p$, $x$ = số lần $A$ xảy ra:
\[
  \frac{x}{n} \xrightarrow{P} p \quad (n \to \infty).
\]
Cơ sở lý thuyết cho định nghĩa xác suất thống kê ở Phần đầu: ``khi $n \to \infty$ thì $x/n \to p$''.
\end{dinhly}

\begin{vidu}[title={Phỏng vấn cử tri}]
$X_i \sim \mathrm{Bern}(p)$, $K_n = \frac{1}{n}\sum X_i$. Chebyshev:
$P(|K_n - p| < \varepsilon) \ge 1 - \dfrac{p(1-p)}{n\varepsilon^2} \ge 1 - \dfrac{1}{4n\varepsilon^2}$.

Ví dụ: $\varepsilon = 0.1$, $n = 100$ $\Rightarrow$ sai số $> 0.1$ với xác suất $\le 0.25$. Muốn tin cậy $95\%$, $\varepsilon = 0.01$ cần $n$ đủ lớn (giải từ bất đẳng thức).
\end{vidu}

\begin{tongket}[title={Tổng kết}]
\begin{enumerate}
  \item Markov, Chebyshev: giới hạn xác suất lệch khỏi kỳ vọng.
  \item LSL: trung bình mẫu hội tụ tới $EX$ (Chebyshev) hoặc $p$ (Bernoulli).
  \item Xấp xỉ nhị thức $\to$ chuẩn (Moivre--Laplace) là bước đệm trước CLT đầy đủ trong thống kê.
\end{enumerate}
\end{tongket}

\end{document}